% Extensions to spectral-braiding.tex
% Five sections: classical r-matrix, pronilpotent Yangian, higher-genus braiding,
% Hopf structure and representation theory, Fredholm partition functions and sewing.

%%%%%%%%%%%%%%%%%%%%%%%%%%%%%%%%%%%%%%%%%%%%%%%%%%%%%%%%%%%%%%%%%%%%%%%%%%%%%
% THE CLASSICAL r-MATRIX AND ITS EXPLICIT FORM
%%%%%%%%%%%%%%%%%%%%%%%%%%%%%%%%%%%%%%%%%%%%%%%%%%%%%%%%%%%%%%%%%%%%%%%%%%%%%

\section{The classical $r$-matrix and its explicit form}
% label removed: sec:classical-r-matrix-explicit
\index{classical r-matrix!explicit form|textbf}
\index{collision residue!r-matrix derivation}

The classical limit of the spectral $R$-matrix was introduced in
\S\ref{subsec:classical-r} as a formal consequence of the quantum
Yang--Baxter equation.  This section extracts the explicit $r$-matrix
for each standard family from the collision residue integral
$\Res^{\mathrm{coll}}_{0,2}(\Theta_\cA)$, building the bridge between
the universal MC element of Volume~I and the concrete spectral data of
integrable systems.

\subsection{The classical limit}
% label removed: subsec:classical-limit-explicit

\begin{definition}[Classical $r$-matrix via semiclassical expansion]
% label removed: def:classical-r-matrix-expansion
\index{classical r-matrix!definition}
Let $R(z) \in \End(L_1 \otimes L_2)[[z^{-1}, z]]$ be the spectral
braiding of Definition~\ref{def:spectral-braiding}, depending on the
loop parameter $\hh$.  The \emph{classical $r$-matrix} is
\begin{equation}
% label removed: eq:classical-r-def
r(z) \;:=\; \lim_{\hh \to 0} \frac{R(z) - \id}{\hh}
\;\in\; \End(L_1 \otimes L_2)((z^{-1})).
\end{equation}
The limit is taken in the $\hh$-adic topology on
$\End(L_1 \otimes L_2)[[z^{-1}, z]][[\hh]]$.  The classical
$r$-matrix is the infinitesimal generator of the spectral braiding:
$R(z) = \id + \hh\, r(z) + O(\hh^2)$.
\end{definition}

\begin{remark}[Collision residue identification]
% label removed: rem:collision-residue-identification
\index{collision residue!identification with r-matrix}
The classical $r$-matrix is the genus-$0$, arity-$2$ projection of
the universal Maurer--Cartan element $\Theta_\cA \in
\MC(\gAmod)$ (Volume~I, Theorem~\ref*{V1-thm:mc2-bar-intrinsic}):
\begin{equation}
% label removed: eq:collision-residue-r
r(z) \;=\; \Res^{\mathrm{coll}}_{0,2}(\Theta_\cA),
\end{equation}
where $\Res^{\mathrm{coll}}_{0,2}$ extracts the genus-$0$,
arity-$2$ collision residue.  This is the binary shadow of the
full shadow obstruction tower: the simplest nontrivial projection
of $\Theta_\cA$, carrying the Lie-algebraic structure of the
theory at the level of pairs of operators.
\end{remark}

\subsection{Derivation of the classical Yang--Baxter equation}
% label removed: subsec:cybe-derivation

The CYBE was stated in \S\ref{subsec:classical-r}; we now give the
complete derivation from the quantum YBE, tracking all signs and
making the algebraic mechanism fully explicit.

\begin{theorem}[Classical Yang--Baxter equation from quantum YBE;
\ClaimStatusProvedHere]
% label removed: thm:cybe-from-qybe
\index{classical Yang--Baxter equation!derivation from quantum}
Let $R(z) = \id + \hh\, r(z) + \hh^2\, s(z) + O(\hh^3)$ satisfy the
quantum Yang--Baxter equation
\begin{equation}
% label removed: eq:qybe-reminder
R_{12}(z_{12})\, R_{13}(z_{13})\, R_{23}(z_{23})
\;=\;
R_{23}(z_{23})\, R_{13}(z_{13})\, R_{12}(z_{12}).
\end{equation}
Then $r(z)$ satisfies the classical Yang--Baxter equation
\begin{equation}
% label removed: eq:cybe-proved
\boxed{
[r_{12}(z_{12}),\, r_{13}(z_{13})]
+ [r_{12}(z_{12}),\, r_{23}(z_{23})]
+ [r_{13}(z_{13}),\, r_{23}(z_{23})]
\;=\; 0.
}
\end{equation}
The order-$\hh^2$ coefficient $s(z)$ satisfies the
\emph{linearized} Yang--Baxter equation with inhomogeneity determined
by $r(z)$.
\end{theorem}

\begin{proof}
\noindent\textbf{Step 1: Expansion of each factor.}
We expand each $R$-matrix factor in~\eqref{eq:qybe-reminder} to
order $\hh^2$:
\begin{align}
R_{12}(z_{12}) &= \id_{123} + \hh\, r_{12} + \hh^2\, s_{12}
+ O(\hh^3), % label removed: eq:R12-expand \\
R_{13}(z_{13}) &= \id_{123} + \hh\, r_{13} + \hh^2\, s_{13}
+ O(\hh^3), % label removed: eq:R13-expand \\
R_{23}(z_{23}) &= \id_{123} + \hh\, r_{23} + \hh^2\, s_{23}
+ O(\hh^3), % label removed: eq:R23-expand
\end{align}
where we suppress the spectral arguments for readability:
$r_{12} \equiv r_{12}(z_{12})$, etc.

\medskip
\noindent\textbf{Step 2: Left-hand side to order $\hh^2$.}
The LHS of~\eqref{eq:qybe-reminder} is the ordered product
$R_{12} R_{13} R_{23}$.  Expanding:
\begin{align*}
\text{LHS} &= (\id + \hh\, r_{12} + \hh^2\, s_{12})
(\id + \hh\, r_{13} + \hh^2\, s_{13})
(\id + \hh\, r_{23} + \hh^2\, s_{23}) + O(\hh^3) \\
&= \id + \hh(r_{12} + r_{13} + r_{23}) \\
&\quad + \hh^2(s_{12} + s_{13} + s_{23}
+ r_{12}\, r_{13} + r_{12}\, r_{23} + r_{13}\, r_{23})
+ O(\hh^3).
\end{align*}
The order-$\hh^2$ terms involving products of $r$'s are:
$r_{12}\, r_{13}$, $r_{12}\, r_{23}$, and $r_{13}\, r_{23}$.

\medskip
\noindent\textbf{Step 3: Right-hand side to order $\hh^2$.}
The RHS $R_{23} R_{13} R_{12}$ expands analogously:
\begin{align*}
\text{RHS} &= \id + \hh(r_{23} + r_{13} + r_{12}) \\
&\quad + \hh^2(s_{23} + s_{13} + s_{12}
+ r_{23}\, r_{13} + r_{23}\, r_{12} + r_{13}\, r_{12})
+ O(\hh^3).
\end{align*}

\medskip
\noindent\textbf{Step 4: Order-$\hh^1$ consistency.}
At order $\hh^1$, both sides give $r_{12} + r_{13} + r_{23}$,
which is consistent: no constraint arises at first order.

\medskip
\noindent\textbf{Step 5: Order-$\hh^2$ identity.}
Setting $\text{LHS} = \text{RHS}$ at order $\hh^2$:
\begin{align*}
&s_{12} + s_{13} + s_{23}
+ r_{12}\, r_{13} + r_{12}\, r_{23} + r_{13}\, r_{23} \\
&\qquad = s_{12} + s_{13} + s_{23}
+ r_{23}\, r_{13} + r_{23}\, r_{12} + r_{13}\, r_{12}.
\end{align*}
The $s$-terms cancel.  Rearranging the remaining terms:
\begin{align*}
r_{12}\, r_{13} - r_{13}\, r_{12}
+ r_{12}\, r_{23} - r_{23}\, r_{12}
+ r_{13}\, r_{23} - r_{23}\, r_{13} &= 0,
\end{align*}
which is
\[
[r_{12},\, r_{13}] + [r_{12},\, r_{23}] + [r_{13},\, r_{23}] = 0.
\]
This is~\eqref{eq:cybe-proved}.

\medskip
\noindent\textbf{Step 6: Linearized equation at order $\hh^3$.}
At order $\hh^3$, the YBE gives
\begin{align*}
&[s_{12}, r_{13}] + [s_{12}, r_{23}]
+ [r_{12}, s_{13}] + [s_{13}, r_{23}]
+ [r_{12}, s_{23}] + [r_{13}, s_{23}] \\
&\qquad + r_{12}\, r_{13}\, r_{23} - r_{23}\, r_{13}\, r_{12} = 0.
\end{align*}
This is the linearized CYBE with inhomogeneity from the cubic
product $r_{12}\, r_{13}\, r_{23} - r_{23}\, r_{13}\, r_{12}$:
the equation governing the first quantum correction $s(z)$ to
the classical $r$-matrix.  The obstruction to deforming a classical
$r$-matrix to a quantum $R$-matrix lives in the third cohomology
of this linearized complex.
\end{proof}

\begin{remark}[The CYBE from Arnold cancellation on $\FM_3(\C)$]
% label removed: rem:cybe-arnold
\index{Arnold cancellation!CYBE}
There is a second, purely geometric proof of the CYBE: apply
Stokes' theorem on $\FM_3(\C)$ to the closed logarithmic
$2$-form
\[
\omega_{3}^{\cl}
\;=\;
\sum_{\{i,j\}} r_0^{(ij)}\,
d\log(z_i - z_j) \wedge d\log(z_k - z_j),
\]
where $\{i,j,k\} = \{1,2,3\}$.  The form $\omega_3^{\cl}$ is
closed on $\Conf_3(\C)$ by the Arnold relation
$d\log(z_{12}) \wedge d\log(z_{23})
+ d\log(z_{23}) \wedge d\log(z_{13})
+ d\log(z_{13}) \wedge d\log(z_{12}) = 0$
and the CYBE for $r_0$.  Conversely, the vanishing of
$\int_{\partial \FM_3(\C)} \omega_3^{\cl}$ forces the CYBE.  Each
boundary face $D_{\{i,j\}}$ of $\FM_3(\C)$ contributes the
commutator $[r_0^{(ij)}, r_0^{(ik)} + r_0^{(jk)}]$; the sum
over the three faces is the CYBE.  The full-collision face
$D_{\{1,2,3\}}$ contributes zero by degree.  This is the
classical limit of the proof of Theorem~\ref{thm:YBE}: the
quantum YBE from $W$-coherence on $\FM_3(\C)$ descends to the
CYBE from the Arnold relation on $\FM_3(\C)$.
\end{remark}

\begin{lemma}[Arnold relation on $\Conf_3(\C)$]
% label removed: lem:arnold-relation
\index{Arnold relation!proof}
Define the logarithmic $1$-forms $\omega_{ij} = d\log(z_i - z_j)$
on $\Conf_3(\C)$.  Then
\begin{equation}% label removed: eq:arnold-relation
\omega_{12} \wedge \omega_{13}
+ \omega_{23} \wedge \omega_{21}
+ \omega_{31} \wedge \omega_{32}
\;=\; 0
\qquad \text{on } \Conf_3(\C).
\end{equation}
\end{lemma}

\begin{proof}
Set $u = z_1 - z_2$, $v = z_2 - z_3$, so that
$z_1 - z_3 = u + v$ and $\omega_{12} = du/u$,
$\omega_{13} = d(u+v)/(u+v)$, $\omega_{23} = dv/v$.
Since $\omega_{ji} = \omega_{ij}$ (the differential
$d\log(-1) = 0$ vanishes), the three terms of~\eqref{eq:arnold-relation} are:
\begin{align*}
\omega_{12} \wedge \omega_{13}
&= \frac{du \wedge dv}{u(u+v)}, &
\omega_{23} \wedge \omega_{21}
&= -\frac{du \wedge dv}{uv}, &
\omega_{31} \wedge \omega_{32}
&= \frac{du \wedge dv}{(u+v)v}.
\end{align*}
Factoring out $du \wedge dv$ and taking the common denominator
$uv(u+v)$:
\[
\frac{v - (u+v) + u}{uv(u+v)}
\;=\;
\frac{0}{uv(u+v)}
\;=\; 0.
\qedhere
\]
\end{proof}

\subsection{Antisymmetry and unitarity}
% label removed: subsec:antisymmetry-unitarity

\begin{proposition}[Antisymmetry and unitarity of $r(z)$;
\ClaimStatusProvedHere]
% label removed: prop:r-antisymmetry-unitarity
\index{classical r-matrix!antisymmetry}
\index{classical r-matrix!unitarity condition}
The classical $r$-matrix satisfies:
\begin{enumerate}[label=\textup{(\roman*)}]
\item \emph{Antisymmetry:}
$r_{12}(z) + r_{21}(-z) = 0$, i.e., $r(-z) = -P \cdot r(z)
\cdot P$ where $P$ is the permutation of tensor factors.
\item \emph{Classical unitarity:}
$r_{12}(z) + r_{12}(-z) = 0$ if $r(z)$ is skew-symmetric
($r_{12} = -r_{21}$), i.e., $r(z)$ is an odd function of~$z$.
\item \emph{Casimir condition:}
If $r(z) = k\,\Omega/z + (\text{regular at } z = 0)$ with
$\Omega = \Omega_{21}$ symmetric (e.g., $\Omega$ is a Casimir)
and $k$ the level \textup{(}AP126\textup{)},
then $r_{12}(z) + r_{21}(-z) = 0$ forces
$r_{\mathrm{reg}}(z) = -r_{\mathrm{reg},21}(-z)$.
\end{enumerate}
\end{proposition}

\begin{proof}
(i) The quantum $R$-matrix satisfies the inverse relation
$R_{12}(z)\, R_{21}(-z) = \id$ (Definition~\ref{def:spectral-braiding}:
composing the braiding $L_1 \otimes L_2 \to L_2 \otimes L_1$
with the reversed braiding gives the identity).  At order $\hh$:
\[
\id + \hh(r_{12}(z) + r_{21}(-z)) + O(\hh^2) = \id,
\]
so $r_{12}(z) + r_{21}(-z) = 0$.

(ii) If $r_{12} = -r_{21}$, then
$r_{12}(z) + r_{21}(-z) = r_{12}(z) - r_{12}(-z) = 0$ gives
$r_{12}(z) = r_{12}(-z)$.  But antisymmetry says $r(-z) = -r_{21}(z)
= r_{12}(z)$, so $r_{12}(-z) = r_{12}(z)$.  Combined: $r(z)$ is
even.  However, for the standard convention where
$r_{21}(z) = P \cdot r(z) \cdot P$, the condition
$r_{12}(z) + r_{21}(-z) = 0$ combined with the simple pole
$k\,\Omega/z$ (level $k$) forces oddness if $\Omega$ is symmetric and
$r - k\,\Omega/z$ is regular.

(iii) Write $r_{12}(z) = k\,\Omega/z + r_{\mathrm{reg}}(z)$
at level $k$ \textup{(}AP126\textup{)}.
Then $r_{21}(-z) = k\,\Omega_{21}/(-z) + r_{\mathrm{reg},21}(-z)
= -k\,\Omega/z + r_{\mathrm{reg},21}(-z)$ (using
$\Omega = \Omega_{21}$).  The antisymmetry condition
$r_{12}(z) + r_{21}(-z) = 0$ becomes
$r_{\mathrm{reg}}(z) + r_{\mathrm{reg},21}(-z) = 0$.
\end{proof}

\subsection{Explicit $r$-matrices for the standard landscape}
% label removed: subsec:explicit-r-matrices

We now compute the classical $r$-matrix for each standard family by
evaluating the collision residue integral
$\Res^{\mathrm{coll}}_{0,2}(\Theta_\cA)$.  In each case, the
computation reduces to extracting the singular part of the
$\lambda$-bracket via the Laplace formula of
Proposition~\ref{prop:field-theory-r}.

\subsubsection{Heisenberg $\cH_k$: the Gaussian $r$-matrix}
% label removed: subsubsec:heisenberg-r-matrix
\index{Heisenberg algebra!classical r-matrix|textbf}
\index{classical r-matrix!Heisenberg}

\begin{proposition}[Heisenberg $r$-matrix from collision residue;
\ClaimStatusProvedHere]
% label removed: prop:heisenberg-r-matrix
For the Heisenberg algebra $\cH_k$ at level $k$, the classical
$r$-matrix is
\begin{equation}
% label removed: eq:heisenberg-r-explicit
\boxed{
r_{\cH_k}(z) \;=\; \frac{k}{z^2}\, J \otimes J,
}
\end{equation}
where $J$ is the Heisenberg current.  This is a scalar $r$-matrix:
it acts diagonally in the one-dimensional current space.
\end{proposition}

\begin{proof}
\noindent\textbf{Step 1: The $\lambda$-bracket.}
The Heisenberg $\lambda$-bracket is
$\{J {}_\lambda J\} = k\lambda$
(Section~\ref{sec:rosetta-stone}).  The coefficients in the
expansion $\{J {}_\lambda J\} = c_0 + c_1 \lambda + c_2 \lambda^2
+ \cdots$ are $c_0 = 0$, $c_1 = k$, $c_n = 0$ for $n \ge 2$.

\medskip
\noindent\textbf{Step 2: The Laplace integral.}
By Proposition~\ref{prop:field-theory-r}, the classical
$r$-matrix is the Laplace transform of the $\lambda$-bracket
kernel:
\[
r(z) = \int_0^\infty d\lambda\; e^{-\lambda z}\,
\{J {}_\lambda J\}\, (J \otimes J)
= \int_0^\infty d\lambda\; e^{-\lambda z}\, k\lambda
\, (J \otimes J).
\]
The integral $\int_0^\infty \lambda\, e^{-\lambda z}\, d\lambda
= 1/z^2$ (for $\Re(z) > 0$) gives
$r(z) = k\, (J \otimes J)/z^2$.

\medskip
\noindent\textbf{Step 3: Laplace kernel and collision residue.}
The formula $r^L(z) = k\, (J \otimes J)/z^2$ is the Laplace
kernel (OPE generating function).  The bar-theoretic collision
residue absorbs the $d\log$ measure, reducing pole orders by
one (cf.\ the Virasoro computation in Step~3 below):
\[
r^{\mathrm{coll}}(z) = \frac{k}{z}\, J \otimes J.
\]
The double pole in the OPE $J(z_1) J(z_2) \sim k/(z_1 - z_2)^2$
becomes a simple pole in $r^{\mathrm{coll}}$, and the absence
of a first-order OPE pole means $r^{\mathrm{coll}}$ has no
regular ($z^0$) term.

\medskip
\noindent\textbf{Step 4: CYBE verification.}
The CYBE~\eqref{eq:cybe-proved} for $r(z) = k\, (J \otimes J)/z^2$
reads
\[
\frac{k^2}{z_{12}^2 z_{13}^2}\, [J_1 J_2, J_1 J_3]
+ \frac{k^2}{z_{12}^2 z_{23}^2}\, [J_1 J_2, J_2 J_3]
+ \frac{k^2}{z_{13}^2 z_{23}^2}\, [J_1 J_3, J_2 J_3] = 0.
\]
Since $J$ is abelian (all commutators $[J_i, J_j] = 0$ as
operators in the classical limit), every commutator vanishes
identically.  The Heisenberg satisfies the CYBE trivially: the
Gaussian $r$-matrix is an abelian solution.
\end{proof}

\begin{remark}[Shadow depth interpretation]
% label removed: rem:heisenberg-r-shadow-depth
The Heisenberg $r$-matrix $r(z) = k/z^2$ is the complete classical
data: no higher corrections arise.  This reflects the shadow depth
$r_{\max}(\cH_k) = 2$ (Gaussian class in the $G/L/C/M$
classification of Volume~I): the shadow obstruction tower terminates at arity~$2$,
and the full Maurer--Cartan element $\Theta_{\cH_k}$ is exhausted
by its binary projection $r(z)$.
\end{remark}

\subsubsection{Affine $\widehat{\fg}_k$: the Casimir $r$-matrix}
% label removed: subsubsec:affine-r-matrix
\index{affine Kac--Moody!classical r-matrix|textbf}
\index{classical r-matrix!affine}
\index{Casimir element!in r-matrix}

\begin{theorem}[Affine $r$-matrix from collision residue;
\ClaimStatusProvedHere]
% label removed: thm:affine-r-matrix
For the affine Kac--Moody algebra $\widehat{\fg}_k$ at level $k$
(with $k + h^\vee \ne 0$), the classical $r$-matrix is
\begin{equation}
% label removed: eq:affine-r-explicit
\boxed{
r_{\widehat{\fg}_k}(z)
\;=\;
\frac{k\,\Omega}{(k + h^\vee)\, z}
\;+\;
\frac{k\, \kappa}{z^2},
}
\end{equation}
where $\Omega = \sum_a t^a \otimes t_a \in \fg \otimes \fg$ is the
quadratic Casimir (with $t^a$ dual to $t_a$ under the Killing form
$\kappa$), and $h^\vee$ is the dual Coxeter number.  The simple pole
$k\,\Omega/((k+h^\vee)z)$ carries the Lie-algebraic data; the double
pole $k\kappa/z^2$ is the Heisenberg-type (abelian Cartan)
contribution. (AP126: both terms carry the level prefix $k$; at
$k=0$ the $r$-matrix vanishes identically.)
\end{theorem}

\begin{proof}
\noindent\textbf{Step 1: The $\lambda$-bracket.}
The affine $\lambda$-bracket is
$\{J^a {}_\lambda J^b\} = f^{ab}_c\, J^c + k\, \kappa^{ab}\, \lambda$,
where $f^{ab}_c$ are the structure constants and $\kappa^{ab}$ is the
invariant bilinear form normalized so that long roots have squared
length~$2$.  The coefficients are $c_0^{ab} = f^{ab}_c\, J^c$ and
$c_1^{ab} = k\, \kappa^{ab}$.

\medskip
\noindent\textbf{Step 2: The Laplace integral.}
The $r$-matrix is
\begin{align*}
r(z) &= \int_0^\infty d\lambda\; e^{-\lambda z}\,
\{J^a {}_\lambda J^b\}\, (t_a \otimes t_b) \\
&= \int_0^\infty d\lambda\; e^{-\lambda z}\,
(f^{ab}_c\, J^c + k\,\kappa^{ab}\,\lambda)\,
(t_a \otimes t_b) \\
&= \frac{f^{ab}_c\, J^c}{z}\, (t_a \otimes t_b)
+ \frac{k\,\kappa^{ab}}{z^2}\, (t_a \otimes t_b).
\end{align*}
The first term is $(\sum_a t^a \otimes [t_a, -])/z$ acting on the
current $J$, which by the Sugawara construction equals
$\Omega/((k + h^\vee)z)$ at the classical level (the factor
$k + h^\vee$ arises from the Sugawara normalization of the
energy-momentum tensor).

Explicitly: the coefficient $f^{ab}_c\, J^c\, (t_a \otimes t_b)/z$
is the adjoint action of the Casimir evaluated on the current.  In
tensor notation, $\sum_{a,b} f^{ab}_c\, t_a \otimes t_b = [\Omega,
t_c \otimes 1 + 1 \otimes t_c]/(k + h^\vee)$.  The identification
with $\Omega/((k + h^\vee)z)$ follows from the defining property
of the Casimir: $\sum_a [t^a, x] \otimes t_a + \sum_a t^a \otimes
[t_a, x] = 0$ for all $x \in \fg$, which is the intertwining
property of $\Omega$.

The second term $k\,\kappa^{ab}\, (t_a \otimes t_b)/z^2 = k\,
\kappa/z^2$ is the abelian (Cartan) contribution, identical in
structure to the Heisenberg $r$-matrix.

\medskip
\noindent\textbf{Step 3: Laplace kernel and collision residue.}
On $\FM_2(\C)$ with collision coordinate $w = z_1 - z_2$, the
genus-$0$ arity-$2$ component of $\Theta_{\widehat{\fg}_k}$ is
\[
\Theta^{(0,2)}(z_1, z_2)
= \frac{f^{ab}_c\, J^c(z_2)}{w}\, t_a \otimes t_b
+ \frac{k\,\kappa^{ab}}{w^2}\, t_a \otimes t_b
+ O(w^0).
\]
The Laplace kernel~\eqref{eq:affine-r-explicit} reproduces these
OPE poles.  The bar-theoretic collision residue absorbs the
$d\log$ measure, reducing pole orders by one: the double-pole
term $k\kappa^{ab}/w^2$ becomes a simple pole in
$r^{\mathrm{coll}}$, and the simple-pole structure-constant term
becomes a regular contribution.

\medskip
\noindent\textbf{Step 4: CYBE verification.}
We verify~\eqref{eq:cybe-proved} for $r(z) = k\,\Omega/z$
at level $k$ \textup{(}AP126: the level survives $d$-log absorption;
at $k = 0$ the $r$-matrix vanishes identically, as required\textup{)},
suppressing the overall $(k+h^\vee)^{-1}$ and the double-pole term, which
satisfies the CYBE independently by the abelian argument of
Proposition~\ref{V1-prop:heisenberg-r-matrix}.

For $r_{ij} = k\,\Omega_{ij}/z_{ij}$, the CYBE reads
\[
k^2\left(
\frac{[\Omega_{12}, \Omega_{13}]}{z_{12}\, z_{13}}
+ \frac{[\Omega_{12}, \Omega_{23}]}{z_{12}\, z_{23}}
+ \frac{[\Omega_{13}, \Omega_{23}]}{z_{13}\, z_{23}}
\right)
\;=\; 0.
\]
The level factor $k^2$ pulls out of every bracket, so the
structural identity reduces to the level-independent
infinitesimal braid relation below; at $k = 0$ both sides
vanish trivially.
The Lie-algebraic factor vanishes
independently of the spectral parameters.  It suffices to
prove the \emph{infinitesimal braid relation}
(Drinfeld--Kohno relation):
\begin{equation}% label removed: eq:infinitesimal-braid
[\Omega_{12}, \Omega_{13}]
+ [\Omega_{12}, \Omega_{23}]
+ [\Omega_{13}, \Omega_{23}]
\;=\; 0
\qquad \text{in } U(\fg)^{\otimes 3}.
\end{equation}
This holds for $\Omega = \sum_a e_a \otimes e^a$ the quadratic
Casimir of any semisimple $\fg$.  Write
$\Omega_{12} = \sum_a e_a^{(1)}\, e^{a\,(2)}$,
$\Omega_{13} = \sum_b e_b^{(1)}\, e^{b\,(3)}$,
$\Omega_{23} = \sum_c e_c^{(2)}\, e^{c\,(3)}$,
where superscripts denote tensor slots.  Then:
\begin{align*}
[\Omega_{12}, \Omega_{13}]
&= \sum_{a,b} [e_a, e_b]^{(1)} \otimes e^{a\,(2)} \otimes e^{b\,(3)}
= \sum_{a,b,d} f_{ab}^d\, e_d^{(1)} \otimes e^{a\,(2)}
  \otimes e^{b\,(3)}, \\
[\Omega_{12}, \Omega_{23}]
&= \sum_{a,c} e_a^{(1)} \otimes [e^a, e_c]^{(2)} \otimes e^{c\,(3)}
= \sum_{a,c,d} f^{ad}_{\phantom{ad}c}\, e_a^{(1)} \otimes e_d^{(2)}
  \otimes e^{c\,(3)}, \\
[\Omega_{13}, \Omega_{23}]
&= \sum_{b,c} e_b^{(1)} \otimes e_c^{(2)} \otimes [e^b, e^c]^{(3)}
= \sum_{b,c,d} f^{bc}_{\phantom{bc}d}\, e_b^{(1)} \otimes e_c^{(2)}
  \otimes e^{d\,(3)}.
\end{align*}
Re-indexing all three sums to a common triple $(p,q,r)$ of
tensor-slot indices, the total coefficient of
$e_p^{(1)} \otimes e_q^{(2)} \otimes e_r^{(3)}$ is
$f_{pr}^q + f^{pq}_r + f^{qr}_p$, which vanishes by the
Jacobi identity $f_{ab}^c + f_{ca}^b + f_{bc}^a = 0$.
Hence~\eqref{eq:infinitesimal-braid} holds, and the CYBE
for $r(z) = k\,\Omega/z$ at level $k$ follows: the three terms in the CYBE
share the common Lie-algebraic factor
$[\Omega_{12}, \Omega_{13}] + [\Omega_{12}, \Omega_{23}]
+ [\Omega_{13}, \Omega_{23}] = 0$,
so their sum vanishes regardless of the spectral denominators
$1/(z_{ij}z_{ik})$ and the overall $k^2$ level factor.
\end{proof}

\begin{remark}[The classical $r$-matrix for $\widehat{\mathfrak{sl}_2}_k$]
% label removed: rem:sl2-r-matrix-explicit
\index{classical r-matrix!$\mathfrak{sl}_2$}
For $\fg = \mathfrak{sl}_2$ with standard generators $e, f, h$
and Killing form $\kappa(h,h) = 2$, $\kappa(e,f) = 1$, the Casimir
is $\Omega = h \otimes h/2 + e \otimes f + f \otimes e$.  The
$r$-matrix is
\[
r_{\widehat{\mathfrak{sl}_2}_k}(z)
= \frac{1}{(k+2)z}
\left(\frac{h \otimes h}{2} + e \otimes f + f \otimes e\right)
+ \frac{k}{z^2}\,\kappa,
\]
where $h^\vee = 2$ for $\mathfrak{sl}_2$.  This is the rational
$r$-matrix of the Yangian $Y(\mathfrak{sl}_2)$, recovered from the
collision residue of the OPE
$J^a(z_1)\, J^b(z_2) \sim f^{ab}_c\, J^c/(z_1 - z_2) +
k\,\kappa^{ab}/(z_1 - z_2)^2$.
\end{remark}

\subsubsection{Virasoro $\mathrm{Vir}_c$: the quartic $r$-matrix}
% label removed: subsubsec:virasoro-r-matrix
\index{Virasoro algebra!classical r-matrix|textbf}
\index{classical r-matrix!Virasoro}

\begin{theorem}[Virasoro OPE kernel and collision residue;
\ClaimStatusProvedHere]
% label removed: thm:virasoro-r-matrix
For the Virasoro algebra $\mathrm{Vir}_c$ at central charge $c$,
the Laplace kernel (OPE generating function) is
\begin{equation}
% label removed: eq:virasoro-r-explicit
\boxed{
r^L_{\mathrm{Vir}_c}(z)
\;=\;
\frac{c}{2\,z^4}
\;+\;
\frac{2\,T}{z^2}
\;+\;
\frac{\partial T}{z},
}
\end{equation}
where $T = T \otimes 1 + 1 \otimes T$ acts on the tensor product
via the energy-momentum tensor, and the identity in $\fg \otimes \fg$
is understood (for each pole, the tensor structure distributes the
action across the two copies of $\mathrm{Vir}_c$).  More precisely,
\begin{equation}
% label removed: eq:virasoro-r-tensor
r^L_{\mathrm{Vir}_c}(z)
= \frac{c/2}{z^4}\, \mathbf{1} \otimes \mathbf{1}
+ \frac{2}{z^2}\, T \otimes \mathbf{1}
+ \frac{1}{z}\, \partial T \otimes \mathbf{1}
+ (\text{regular at } z = 0),
\end{equation}
where we write the ``first slot acts'' convention.
The collision residue
$r^{\mathrm{coll}}_{\mathrm{Vir}_c}(z)
= \operatorname{Res}^{\mathrm{coll}}_{0,2}(\Theta_{\mathrm{Vir}_c})$
has pole orders one lower, since the $d\log$ kernel on
$\FM_2(\C)$ absorbs one power of~$z$
\textup{(}Volume~I, Remark~\textup{\ref{V1-rem:three-r-matrices}}\textup{)}:
\begin{equation}
r^{\mathrm{coll}}_{\mathrm{Vir}_c}(z)
\;=\;
\frac{c/2}{z^3}\,\mathbf{1}\otimes\mathbf{1}
\;+\;
\frac{2T}{z}\,\otimes\,\mathbf{1}.
\end{equation}
\end{theorem}

\begin{proof}
\noindent\textbf{Step 1: The $\lambda$-bracket.}
The Virasoro $\lambda$-bracket is
$\{T {}_\lambda T\} = (\partial + 2\lambda)\, T +
\frac{c}{12}\,\lambda^3$.
Expanding: $c_0 = \partial T$, $c_1 = 2T$, $c_2 = 0$,
$c_3 = c/12$.

\medskip
\noindent\textbf{Step 2: The Laplace integral.}
\begin{align*}
r^L(z) &= \int_0^\infty d\lambda\; e^{-\lambda z}\,
\{T {}_\lambda T\} \\
&= \int_0^\infty d\lambda\; e^{-\lambda z}\,
\bigl(\partial T + 2\lambda\, T + \tfrac{c}{12}\,\lambda^3\bigr) \\
&= \frac{\partial T}{z} + \frac{2T}{z^2}
+ \frac{c}{12} \cdot \frac{3!}{z^4} \\
&= \frac{\partial T}{z} + \frac{2T}{z^2} + \frac{c}{2z^4},
\end{align*}
using $\int_0^\infty \lambda^n\, e^{-\lambda z}\, d\lambda
= n!/z^{n+1}$.

\medskip
\noindent\textbf{Step 3: From Laplace kernel to collision residue.}
The Laplace kernel $r^L(z)$ is the OPE generating function: it
reproduces the position-space OPE
$T(z_1)\, T(z_2) \sim c/(2w^4) + 2T(z_2)/w^2 + \partial T(z_2)/w$
(where $w = z_1 - z_2$).
The bar-theoretic collision residue
$r^{\mathrm{coll}}(z) =
\Res^{\mathrm{coll}}_{0,2}(\Theta_{\mathrm{Vir}_c})$
extracts residues along $d\log(z_1 - z_2)$ rather than
$dz/(z_1 - z_2)$; the $d\log$ measure absorbs one power of~$w$,
so the pole orders in $r^{\mathrm{coll}}$ are one lower than
in the OPE:
$r^{\mathrm{coll}}(z) = (c/2)/z^3 + 2T/z$.
In particular, for bosonic algebras the collision residue has no
even-order poles ($z^{-2n}$ in the OPE maps to $z^{-(2n-1)}$ in
$r^{\mathrm{coll}}$).

\medskip
\noindent\textbf{Step 4: Pole structure and shadow depth.}
The quartic pole in $r^L(z)$ (equivalently, the cubic pole in
$r^{\mathrm{coll}}(z)$) is the signature of the Virasoro: no
simple Lie algebra produces a pole of order higher than~$2$ in the
Laplace kernel (the affine $r$-matrix has at most a double pole).
The anomalous pole arises from the $\lambda^3$ term in the
$\lambda$-bracket, which encodes the failure of the
energy-momentum tensor to transform as a primary field.
In the shadow depth classification, this higher-order pole places
the Virasoro in the $M$ (mixed) class: the shadow obstruction tower does not
terminate, and the full Maurer--Cartan element
$\Theta_{\mathrm{Vir}_c}$ requires infinitely many arity
corrections beyond the binary $r$-matrix.
\end{proof}

\subsubsection{$\mathcal{W}_3$: the sextic $r$-matrix}
% label removed: subsubsec:w3-r-matrix
\index{W-algebra!$\mathcal{W}_3$ r-matrix|textbf}
\index{classical r-matrix!$\mathcal{W}_3$}

\begin{proposition}[$\mathcal{W}_3$ $r$-matrix;
\ClaimStatusProvedHere]
% label removed: prop:w3-r-matrix
For the $\mathcal{W}_3$ algebra at central charge $c$
(generated by the Virasoro field $T$ of weight~$2$ and the
spin-$3$ current $W$ of weight~$3$), the classical $r$-matrix
has the form
\begin{equation}
% label removed: eq:w3-r-explicit
r_{\mathcal{W}_3}(z)
\;=\;
r_{\mathrm{Vir}_c}(z)
\;+\;
\frac{W_0}{z^3}
\;+\;
\frac{W_1}{z^4}
\;+\;
\frac{W_2}{z^5}
\;+\;
\frac{c_W}{z^6},
\end{equation}
where $r_{\mathrm{Vir}_c}(z)$ is the Virasoro contribution
\eqref{eq:virasoro-r-explicit}, $W_0, W_1, W_2$ are
$\lambda$-bracket coefficients of the $W$-current OPE, and
$c_W$ is the $W$-algebra central term.  The sextic pole
$c_W/z^6$ arises from the OPE $W(z_1)\, W(z_2) \sim
c_W/(z_1 - z_2)^6 + \cdots$ and is the defining signature
of the spin-$3$ field.
\end{proposition}

\begin{proof}
The $\lambda$-bracket structure of $\mathcal{W}_3$ involves three
brackets: $\{T {}_\lambda T\}$ (Virasoro, as above),
$\{T {}_\lambda W\} = (\partial + 3\lambda)\, W$ (primary of
weight~$3$), and $\{W {}_\lambda W\}$ (the $\mathcal{W}_3$ OPE,
involving $T$, $\Lambda = (TT) - (3/10)\partial^2 T$, and a
quintic polynomial in $\lambda$).  The $\lambda$-expansion
$\{W {}_\lambda W\} = \sum_{n=0}^{5} c_n^{WW}\, \lambda^n$
has $c_5^{WW} = c/3\cdot5!$ (the highest pole, contributing
$c_W/z^6$) and lower coefficients involving the composite field
$\Lambda$ and its derivatives.

The Laplace transform $r(z) = \int_0^\infty e^{-\lambda z}\,
\{W {}_\lambda W\}\, d\lambda$ produces a sum of poles
$c_n^{WW}\, n!/z^{n+1}$ for $n = 0, \ldots, 5$, with the
leading term $c_5^{WW}\, 5!/z^6 = c_W/z^6$.  Combined with the
$T{-}T$ and $T{-}W$ brackets, this gives~\eqref{eq:w3-r-explicit}.
\end{proof}

\begin{remark}[General $\mathcal{W}_N$: $r$-matrix pole order]
% label removed: rem:wn-r-matrix-pole
\index{W-algebra!pole order of r-matrix}
For the $\mathcal{W}_N$ algebra (generated by fields of spins
$2, 3, \ldots, N$), the highest pole in the $r$-matrix is
$r(z) \sim c_{W_N}/z^{2N}$, arising from the self-OPE of the
spin-$N$ generator $W_N$.  The pole order $2N$ grows linearly
with $N$: each additional generator contributes higher poles to
the collision residue.  In the limit $N \to \infty$
($\mathcal{W}_\infty$), the $r$-matrix has essential singularity
at $z = 0$; the formal Laurent series does not terminate.
\end{remark}

\subsubsection{Beta-gamma $\beta\gamma_\lambda$: the fermionic
$r$-matrix}
% label removed: subsubsec:betagamma-r-matrix
\index{beta-gamma system!classical r-matrix|textbf}
\index{classical r-matrix!beta-gamma}

\begin{proposition}[Beta-gamma $r$-matrix from collision residue;
\ClaimStatusProvedHere]
% label removed: prop:betagamma-r-matrix
For the beta-gamma system $\beta\gamma_\lambda$ at conformal weight
$\lambda$ (with $\beta$ of weight $\lambda$ and $\gamma$ of weight
$1 - \lambda$), the classical $r$-matrix is
\begin{equation}
% label removed: eq:betagamma-r-explicit
r_{\beta\gamma_\lambda}(z)
\;=\;
\frac{1}{z}\, (\beta \otimes \gamma + \gamma \otimes \beta),
\end{equation}
with a single simple pole.  The $r$-matrix is \emph{non-diagonal}
in field space: it exchanges $\beta$ and $\gamma$ across the two
tensor factors.
\end{proposition}

\begin{proof}
The $\lambda$-bracket of the beta-gamma system is
$\{\beta {}_\lambda \gamma\} = 1$ (contact OPE
$\beta(z_1)\, \gamma(z_2) \sim 1/(z_1 - z_2)$), with
$\{\beta {}_\lambda \beta\} = \{\gamma {}_\lambda \gamma\} = 0$.
The only nonzero coefficient is $c_0^{\beta\gamma} = 1$.

The Laplace transform gives
$r(z) = \int_0^\infty e^{-\lambda z}\, d\lambda\,
(\beta \otimes \gamma + \gamma \otimes \beta)
= (\beta \otimes \gamma + \gamma \otimes \beta)/z$.

The collision residue on $\FM_2(\C)$ is
$\Res_{w=0}\, [1/w\, \beta(z_2) \otimes \gamma(z_2)\, dw]
= \beta \otimes \gamma$, plus the $\gamma \otimes \beta$
term from the opposite ordering.  The simple pole structure
confirms the contact shadow depth
$r_{\max}(\beta\gamma_\lambda) = 4$ arises not from the
binary $r$-matrix but from the \emph{quartic} self-interaction
$\beta\beta\gamma\gamma$ at arity~$4$ (Volume~I,
\S\ref*{V1-sec:betagamma-quartic-birth}).
\end{proof}

\begin{remark}[Comparison of pole structures across the landscape]
% label removed: rem:pole-structure-comparison
\index{classical r-matrix!pole structure comparison}
The pole structure of $r(z)$ at $z = 0$ encodes the conformal
spin content of the chiral algebra:
\begin{center}
\small
\renewcommand{\arraystretch}{1.15}
\begin{tabular}{llll}
\textbf{Family} & \textbf{Highest pole} & \textbf{Source} &
\textbf{Shadow class} \\
\hline
$\cH_k$ & $k/z^2$ [Laplace] & $\lambda$-bracket $k\lambda$ & $G$ \\
$\widehat{\fg}_k$ & $k\,\Omega/z + k\kappa/z^2$ [Laplace] & structure constants $f^{ab}_c$ & $L$ \\
$\beta\gamma_\lambda$ & $1/z$ [Laplace] & contact OPE & $C$ \\
$\mathrm{Vir}_c$ & $c/(2z^4)$ [Laplace] & conformal anomaly $\lambda^3$ & $M$ \\
$\mathcal{W}_N$ & $c_{W_N}/z^{2N}$ [Laplace] & spin-$N$ self-OPE & $M$ \\
\end{tabular}
\end{center}
The pole order determines the leading behavior of the spectral
$R$-matrix near the collision singularity.  The shadow class
($G$/$L$/$C$/$M$ from Volume~I) is determined by the full tower
$\Theta_\cA$, not just the binary projection $r(z)$: the
Heisenberg has a double pole but terminates at arity~$2$
(Gaussian), while the Virasoro has a quartic pole and an infinite
tower (mixed).
\end{remark}


%%%%%%%%%%%%%%%%%%%%%%%%%%%%%%%%%%%%%%%%%%%%%%%%%%%%%%%%%%%%%%%%%%%%%%%%%%%%%
% THE PRONILPOTENT YANGIAN
%%%%%%%%%%%%%%%%%%%%%%%%%%%%%%%%%%%%%%%%%%%%%%%%%%%%%%%%%%%%%%%%%%%%%%%%%%%%%

\section{The pronilpotent Yangian $Y_T^{\mathrm{mod}}$}
% label removed: sec:pronilpotent-yangian
\index{Yangian!pronilpotent|textbf}
\index{pronilpotent Yangian|textbf}
\index{modular convolution algebra!Yangian}

The classical $r$-matrix is the genus-$0$, arity-$2$ projection of
$\Theta_\cA$.  Beyond genus~$0$, the local theorem surface of this
chapter supplies an ambient pronilpotent completion
$Y_T^{\mathrm{mod}}$ of the modular convolution algebra together with
formal candidate higher-genus binary coefficients, but not a proved
modular kernel.  This section constructs that ambient completion and
records the stable-graph Yang--Baxter equation as the formal
Maurer--Cartan criterion that such a kernel would have to satisfy.

\subsection{The modular convolution algebra}
% label removed: subsec:modular-convolution-yangian

\begin{definition}[Modular convolution algebra for the spectral
problem]
% label removed: def:modular-convolution-spectral
\index{modular convolution algebra!spectral definition}
Let $\cA$ be a chirally Koszul logarithmic $\SCchtop$-algebra, with bar
coalgebra $\barBch(\cA)$ and Koszul dual $\cA^!$.  The
\emph{modular convolution Lie algebra for the spectral problem} is
\begin{equation}
% label removed: eq:modular-convolution-spectral
L_T^{\mathrm{mod}}
\;=\;
\prod_{g \ge 0,\, n \ge 1}
\Hom\bigl(
C_\ast(\overline{\mathcal{M}}_{g, n+1})
\otimes (\cA^!)^{\otimes n},\;
\cA^!
\bigr)[1],
\end{equation}
where $\overline{\mathcal{M}}_{g,n+1}$ is the Deligne--Mumford
moduli space of stable curves of genus $g$ with $n+1$ marked
points (one ``output'' point and $n$ ``input'' points).  The
shift $[1]$ is the cohomological shift placing the convolution
elements in degree $1$ (the Maurer--Cartan degree).

Elements of $L_T^{\mathrm{mod}}$ are sequences
$\{\alpha_{g,n}\}_{g \ge 0,\, n \ge 1}$ of linear maps
\[
\alpha_{g,n} \colon
C_\ast(\overline{\mathcal{M}}_{g,n+1})
\otimes (\cA^!)^{\otimes n}
\longrightarrow \cA^![1]
\]
for each $(g,n)$ with $2g - 2 + n > 0$ (stability).
\end{definition}

\begin{definition}[Lie bracket via graph composition]
% label removed: def:lie-bracket-graph-composition
\index{Lie bracket!graph composition}
\index{graph composition!Lie bracket}
The Lie bracket on $L_T^{\mathrm{mod}}$ is defined by \emph{graph
composition with propagator}.  Given
$\alpha = \{\alpha_{g_1, n_1}\}$ and
$\beta = \{\beta_{g_2, n_2}\}$, the bracket
$[\alpha, \beta]_{g, n}$ is defined by summing over all stable
graphs $\Gamma$ of type $(g, n+1)$ with exactly one internal
edge connecting a vertex of type $(g_1, n_1 + 1)$ (decorated by
$\alpha$) to a vertex of type $(g_2, n_2 + 1)$ (decorated by
$\beta$).  The internal edge carries the propagator $r(z)$, the
classical $r$-matrix evaluated at the spectral parameter of the
edge.

Concretely, for a graph $\Gamma$ with one edge $e$ connecting
vertices $v_1$ (type $(g_1, n_1+1)$) and $v_2$
(type $(g_2, n_2+1)$), with $g = g_1 + g_2 + h_1(\Gamma)$
and $n = n_1 + n_2 - 2$ (after identifying the two half-edges
of $e$ with one input each):
\begin{equation}
% label removed: eq:lie-bracket-graph
[\alpha, \beta]_\Gamma
\;=\;
\sum_{\text{gluings}}
\alpha_{g_1, n_1}
\circ_e
r(z_e)
\circ_e
\beta_{g_2, n_2},
\end{equation}
where $\circ_e$ denotes the composition along the edge $e$ using
the $r$-matrix as propagator, and the sum is over all ways of
distributing the $n$ external half-edges between the two vertices.
The full bracket is the sum over all stable graphs with one edge:
\[
[\alpha, \beta]_{g,n}
\;=\;
\sum_{\Gamma \in \StGraph_{g,n+1}^{(1)}}
\frac{1}{|\Aut(\Gamma)|}\,
[\alpha, \beta]_\Gamma.
\]
\end{definition}

\begin{proposition}[Jacobi identity from codimension-$2$ cancellation;
\ClaimStatusProvedHere]
% label removed: prop:jacobi-graph-composition
\index{Jacobi identity!graph composition}
The bracket~\eqref{eq:lie-bracket-graph} satisfies the Jacobi
identity.  The proof is by codimension-$2$ cancellation in
$\overline{\mathcal{M}}_{g,n+1}$: the three terms of the Jacobi
identity
$[[\alpha, \beta], \gamma] + [[\beta, \gamma], \alpha]
+ [[\gamma, \alpha], \beta] = 0$
correspond to the three ways of connecting three vertices by two
edges in a graph with two internal edges, and these are the three
codimension-$2$ boundary strata of $\overline{\mathcal{M}}_{g,n+1}$
arising from two successive degenerations.  The cancellation follows
from $\partial^2 = 0$ on $\overline{\mathcal{M}}_{g,n+1}$: the
codimension-$2$ boundary of a manifold with corners has total
boundary zero.
\end{proposition}

\begin{proof}
The Lie bracket $[\alpha, \beta]$ is computed by \emph{edge gluing}:
one output half-edge of $\alpha$ is joined to one input half-edge
of $\beta$ (or vice versa) to form a single internal edge.  This
is the Lie bracket structure (not operadic insertion): it connects
an output of one element to an input of another via one new edge,
increasing the total internal edge count by exactly one.  If
$\alpha$ has arity $n_1$ and $\beta$ has arity $n_2$, the
bracket $[\alpha, \beta]$ has arity $n_1 + n_2 - 1$, since one
input from each vertex is consumed by the gluing edge.  The
filtration compatibility $[F^{N_1}, F^{N_2}] \subset
F^{N_1 + N_2 - 1}$ is additive in the edge count.

The double bracket $[[\alpha, \beta], \gamma]$ sums over graphs
with two internal edges: the first edge connects $\alpha$ to
$\beta$ (forming $[\alpha, \beta]$), and the second edge connects
the composite $[\alpha, \beta]$ to $\gamma$.  The three cyclic
permutations of the Jacobi identity correspond to the three ways
of choosing which pair is connected by the first edge:
$\alpha$-$\beta$, $\beta$-$\gamma$, or $\gamma$-$\alpha$.

Each such graph $\Gamma^{(2)}$ with two internal edges corresponds
to a codimension-$2$ boundary stratum
$\overline{\mathcal{M}}_{\Gamma^{(2)}}$ of
$\overline{\mathcal{M}}_{g,n+1}$: the curve degenerates along
two nodes simultaneously.  The two edges give two codimension-$1$
degenerations whose composition is codimension~$2$.  The three
terms of the Jacobi identity correspond to the three orderings
of the two degenerations.

By the general boundary-of-boundary identity for manifolds with
corners, $\partial(\partial \overline{\mathcal{M}}_{g,n+1}) = 0$:
the codimension-$2$ strata, oriented by the two normal covectors,
cancel pairwise.  This is the combinatorial identity underlying
$\partial^2 = 0$ on the chain complex
$C_\ast(\overline{\mathcal{M}}_{g,n+1})$.  Translated to the
Lie algebra, it gives the Jacobi identity for the edge-gluing
bracket.

The sign verification uses the Koszul sign rule on the
desuspension shift: each edge contributes a sign $(-1)^{|e|}$
from the degree-$1$ propagator.  The three terms of the Jacobi
identity acquire signs
$(-1)^{|\alpha|(|\beta| + |\gamma|)}$,
$(-1)^{|\beta|(|\gamma| + |\alpha|)}$,
$(-1)^{|\gamma|(|\alpha| + |\beta|)}$,
which is the graded Jacobi identity in degree $1$.
\end{proof}

\subsection{The filtration and pronilpotent completion}
% label removed: subsec:filtration-pronilpotent

\begin{definition}[Arity filtration]
% label removed: def:arity-filtration
\index{arity filtration!on convolution algebra}
The \emph{arity filtration} on $L_T^{\mathrm{mod}}$ is defined by
\begin{equation}
% label removed: eq:arity-filtration
F^N L_T^{\mathrm{mod}}
\;=\;
\prod_{\substack{g \ge 0,\, n \ge 1 \\ n \ge N}}
\Hom\bigl(
C_\ast(\overline{\mathcal{M}}_{g, n+1})
\otimes (\cA^!)^{\otimes n},\;
\cA^!
\bigr)[1].
\end{equation}
This filtration is separated ($\bigcap_N F^N = 0$) and compatible
with the Lie bracket: $[F^{N_1}, F^{N_2}] \subset
F^{N_1 + N_2 - 1}$ (the bracket of an element of arity $n_1$ with
an element of arity $n_2$ produces an element of arity
$n_1 + n_2 - 1$, since one input from each is consumed by the
internal edge).  The quotients $L_T^{\mathrm{mod}} / F^{N+1}$ are
finite-arity truncations.

The codimension of the boundary stratum $\sigma_T$ corresponding to
a stable tree~$T$ in $\overline{\mathcal{M}}_{g,n+1}$ is the
number of \emph{internal edges} $|E_{\mathrm{int}}(T)|$ of~$T$,
not the number of internal vertices minus one.  For binary trees,
$|E_{\mathrm{int}}(T)| = |V_{\mathrm{int}}(T)| - 1$, so the two
counts coincide; for non-binary trees (vertices of valence
$\ge 4$), the edge count is the correct codimension.
\end{definition}

\begin{definition}[Pronilpotent Yangian]
% label removed: def:pronilpotent-yangian
\index{Yangian!pronilpotent completion}
The \emph{pronilpotent Yangian} is the completion of the modular
convolution Lie algebra with respect to the arity filtration:
\begin{equation}
% label removed: eq:pronilpotent-yangian
Y_T^{\mathrm{mod}}
\;:=\;
\varprojlim_{N} \;
L_T^{\mathrm{mod}} / F^{N+1}.
\end{equation}
Elements of $Y_T^{\mathrm{mod}}$ are formal power series in the
arity variable: a sequence $\{\alpha_{g,n}\}_{g \ge 0,\, n \ge 1}$
with no convergence requirement on the tail $n \to \infty$.  The
pronilpotent Lie algebra $Y_T^{\mathrm{mod}}$ carries a natural
Maurer--Cartan set
\begin{equation}
% label removed: eq:mc-set-yangian
\MC(Y_T^{\mathrm{mod}})
\;=\;
\bigl\{\alpha \in Y_T^{\mathrm{mod}}
\;\bigm|\;
d\alpha + \tfrac{1}{2}[\alpha, \alpha] = 0
\bigr\},
\end{equation}
which is well-defined because the bracket is pronilpotent
(the arity-$N$ component of $[\alpha, \alpha]$ depends only on
the components of $\alpha$ in arities $< N$).
\end{definition}

\begin{remark}[Spectral sequence convergence]
% label removed: rem:spectral-sequence-convergence
\index{spectral sequence!convergence of arity filtration}
The spectral sequence associated to the arity filtration
$F^\bullet$ on $Y_T^{\mathrm{mod}}$ converges: at each total degree~$d$, only finitely many filtration
stages contribute: an element of arity~$n$ and genus~$g$ has
total degree $d = 2g - 2 + n$ (from the stability condition
$2g - 2 + n > 0$ and the grading), so the arity $n \le d + 2$
is bounded by the total degree.  Hence $F^{d+3} = 0$ in total
degree~$d$, and the filtration is finite in each degree.  By the
classical convergence theorem for bounded-below exhaustive
filtrations, the spectral sequence converges to $H^\ast(Y_T^{\mathrm{mod}}, d)$
at each total degree.
\end{remark}

\begin{remark}[Genus expansion]
% label removed: rem:genus-expansion-yangian
\index{genus expansion!pronilpotent Yangian}
There is a second filtration on $Y_T^{\mathrm{mod}}$ by genus:
$G^g Y_T^{\mathrm{mod}}$ consists of elements supported in genera
$\ge g$.  This filtration is compatible with the Lie bracket (genus
is additive under gluing) and gives a double filtration by arity
and genus.  The associated graded with respect to the genus
filtration separates the contributions from different genera:
$\gr^G_g\, Y_T^{\mathrm{mod}}$ consists of genus-$g$ graph sums
only.
\end{remark}

\subsection{A candidate higher-genus $R$-matrix package}
% label removed: subsec:pro-mc-element

\begin{definition}[Candidate stable-graph $R$-matrix]
% label removed: def:stable-graph-r-matrix
\index{R-matrix!stable-graph}
\index{stable graph!R-matrix}
The \emph{candidate stable-graph $R$-matrix} is a formal series
\begin{equation}
% label removed: eq:stable-graph-r-matrix
R_T^{\mathrm{mod}}(z; \hh)
\;=\;
\sum_{g \ge 0} \hh^{2g}\, r_{T,g}(z)
\;\in\;
Y_T^{\mathrm{mod}}[[\hh]],
\end{equation}
where $r_{T,g}(z)$ is intended as the genus-$g$ contribution.  The classical
$r$-matrix of \S\ref{sec:classical-r-matrix-explicit} is the
genus-$0$ term: $r_{T,0}(z) = r(z)$.  On the conjectural modular-kernel
surface, the higher coefficients would encode loop corrections to the
spectral braiding.
\end{definition}

\begin{theorem}[Stable-graph YBE as a formal MC criterion;
\ClaimStatusProvedHere]
% label removed: thm:stable-graph-ybe-mc
\index{Yang--Baxter equation!stable-graph}
\index{Maurer--Cartan equation!stable-graph YBE}
Let
$R(z;\hh)=\sum_{g \ge 0}\hh^{2g}r_g(z)\in Y_T^{\mathrm{mod}}[[\hh]]$
be a formal series with genus-$0$ term $r_0(z)=r(z)$.  Then the
stable-graph Yang--Baxter equations for the coefficients $r_g(z)$ are
equivalent to the Maurer--Cartan equation for $R$ in
$Y_T^{\mathrm{mod}}[[\hh]]$:
\begin{equation}
% label removed: eq:stable-graph-ybe-mc
d\, R
+ \frac{1}{2}\bigl[R,\, R\bigr]
= 0
\quad \Longleftrightarrow \quad
\text{the stable-graph YBE holds coefficientwise in genus.}
\end{equation}
\end{theorem}

\begin{proof}
\noindent\textbf{Step 1: Expansion by genus.}
Write $R = \sum_{g \ge 0} \hh^{2g}\, r_g$.
The MC equation $dR + \frac{1}{2}[R, R] = 0$ expands as
\begin{equation}
% label removed: eq:mc-genus-expansion
\sum_{g \ge 0} \hh^{2g}\, d\, r_g
+ \frac{1}{2} \sum_{g_1, g_2 \ge 0}
\hh^{2(g_1 + g_2)}\, [r_{g_1}, r_{g_2}] = 0.
\end{equation}
At each genus $g$, the coefficient of $\hh^{2g}$ gives
\begin{equation}
% label removed: eq:mc-at-genus-g
d\, r_g
+ \frac{1}{2} \sum_{\substack{g_1 + g_2 = g \\ g_1, g_2 \ge 0}}
[r_{g_1}, r_{g_2}] = 0.
\end{equation}

\medskip
\noindent\textbf{Step 2: Genus $0$.}
At $g = 0$: $d\, r_0 + \frac{1}{2}[r_0, r_0] = 0$.  The
bracket $[r_0, r_0]$ at arity $2$ (acting on pairs of elements
of $\cA^!$) is computed by the single graph with one internal
edge connecting two binary vertices.  In the spectral variable,
this is exactly the left-hand side of the CYBE:
$[r_0, r_0]_{12,13} + [r_0, r_0]_{12,23}
+ [r_0, r_0]_{13,23}$, where the three terms correspond to the
three graphs (one for each pair $\{i,j\}$ of external legs
connected through the propagator).  The differential $d\, r_0$
vanishes if $r_0$ is a cocycle (which holds when $r_0$ is
defined at the cohomology level).  Hence the genus-$0$ MC
equation is the CYBE.

\medskip
\noindent\textbf{Step 3: Genus $1$.}
At $g = 1$: $d\, r_1 + [r_0, r_1] = 0$ (the $[r_1, r_0]$
term equals $[r_0, r_1]$ by antisymmetry, giving the factor
$1/2 \cdot 2$).  The bracket $[r_0, r_1]$ is a genus-$1$
graph with one internal edge: a genus-$0$ vertex carrying $r_0$
connected to a genus-$1$ vertex carrying $r_1$.  In addition,
there is a self-loop term from $\frac{1}{2}[r_0, r_0]$ at
genus~$1$: a genus-$0$ graph with one self-loop connecting two
inputs of the same binary vertex.  This self-loop contributes
the sewing operator $\hh\,\Delta_{\mathrm{sew}}$ acting on
$r_0$.  Thus the genus-$1$ coefficient of the MC equation is
exactly the genus-$1$ stable-graph Yang--Baxter equation.

\medskip
\noindent\textbf{Step 4: General $g$.}
At general genus $g$, equation~\eqref{eq:mc-at-genus-g} is the
genus-$g$ stable-graph Yang--Baxter equation: it is the genus-$g$
coefficient of the single formal Maurer--Cartan equation and packages
the lower-genus interaction terms
$\sum_{g_1+g_2=g,\, g_i < g} [r_{g_1}, r_{g_2}]$.

\medskip
\noindent\textbf{Step 5: Equivalence with stable-graph YBE.}
The stable-graph Yang--Baxter equation at genus $g$ and arity $n$
asserts that the sum over all stable graphs $\Gamma$ of type
$(g, n)$ with three external legs (the ``triple interaction'' at
arity $3$) weighted by $R_T^{\mathrm{mod}}$ vanishes.  This is
precisely the arity-$3$ component of the MC
equation~\eqref{eq:mc-at-genus-g}, since the bracket $[r_{g_1},
r_{g_2}]$ at arity $3$ produces the three-term sum indexed by
the three stable graphs with one internal edge and three
external legs.  The equivalence extends to all arities $n$
by the higher MC identities (the $L_\infty$ generalization of
the Jacobi identity, which encodes graph sums at all arities).
\end{proof}

\begin{corollary}[Classical Yangian as genus-$0$ truncation;
\ClaimStatusProvedHere]
% label removed: cor:classical-yangian-truncation
\index{Yangian!classical as genus-0 truncation}
The classical Yangian $Y(\fg)$ is the genus-$0$ truncation of
$Y_T^{\mathrm{mod}}$: the quotient
$Y_T^{\mathrm{mod}} / G^1 Y_T^{\mathrm{mod}}$ retains only the
tree-level (genus-$0$) contribution.  The classical $r$-matrix
$r(z) = r_{T,0}(z)$ is the unique Maurer--Cartan element at this
truncation level.

For $\cA = \widehat{\fg}_k$: the classical Yangian $Y(\fg)$ is
recovered from the genus-$0$ projection of $Y_T^{\mathrm{mod}}$.
The $R$-matrix $R(z) = 1 + \hh\, \Omega/((k+h^\vee)z)
+ O(\hh^2)$ deforms to the quantum Yangian $R$-matrix
$R^{\mathrm{qu}}(z) = (z + \hh\,P)/(z - \hh)$ by resumming the
genus expansion to all orders.
\end{corollary}

\begin{proof}
The genus filtration $G^\bullet$ on $Y_T^{\mathrm{mod}}$ has
$\gr^G_0$ consisting of tree-level graphs only.  The MC equation
at genus $0$ is $d\, r_0 + \frac{1}{2}[r_0, r_0] = 0$, the
CYBE, which determines $r_0 = r(z)$.  The classical Yangian
$Y(\fg)$ is defined as the universal envelope of the relation
$[x, y] = r_0(x \otimes y - y \otimes x)$; this is the content
of the genus-$0$ MC equation applied to pairs of generators.
\end{proof}

\subsection{The gravitational Yangian for Virasoro}
% label removed: subsec:gravitational-yangian-virasoro
\index{Yangian!gravitational|textbf}
\index{Virasoro algebra!gravitational Yangian}

\begin{definition}[Gravitational Yangian]
% label removed: def:gravitational-yangian
\index{gravitational Yangian!definition}
For $\cA = \mathrm{Vir}_c$, the pronilpotent Yangian
$Y_T^{\mathrm{mod}}(\mathrm{Vir}_c)$ is the \emph{gravitational
Yangian}.  Its modular convolution algebra is
\[
L_T^{\mathrm{mod}}(\mathrm{Vir}_c)
= \prod_{g,n}
\Hom\bigl(
C_\ast(\overline{\mathcal{M}}_{g,n+1})
\otimes (\mathrm{Vir}_{26-c})^{\otimes n},\;
\mathrm{Vir}_{26-c}
\bigr)[1],
\]
where $\mathrm{Vir}_{26-c} = \mathrm{Vir}_c^!$ is the Koszul
dual Virasoro at the complementary central charge
(Volume~I, critical pitfall: $\mathrm{Vir}_c^! = \mathrm{Vir}_{26-c}$,
self-dual at $c = 13$, NOT $c = 26$).
\end{definition}

\begin{proposition}[First properties of the gravitational Yangian;
\ClaimStatusProvedHere]
% label removed: prop:gravitational-yangian-properties
The gravitational Yangian
$Y_T^{\mathrm{mod}}(\mathrm{Vir}_c)$ has the following properties:
\begin{enumerate}[label=\textup{(\roman*)}]
\item \emph{Genus-$0$ Laplace kernel:}
$r^L_0(z) = c/(2z^4) + 2T/z^2 + \partial T/z$;
equivalently, collision residue
$r^{\mathrm{coll}}_0(z) = (c/2)/z^3 + 2T/z$
(Theorem~\ref{thm:virasoro-r-matrix}).

\item \emph{Non-formality at all genera:}
The higher operations $m_k \ne 0$ for $k \ge 3$ in the $\Ainf$
structure on $\mathrm{Vir}_{26-c}$ produce nontrivial higher
brackets $\ell_k$ in the $L_\infty$ structure on the convolution
algebra.  The shadow obstruction tower $\Theta_{\mathrm{Vir}_c}^{(\le r)}$
does not terminate: $r_{\max} = \infty$.

\item \emph{Genus-$1$ correction:}
The first quantum correction $r_1(z)$ to the $r$-matrix encodes
the effect of the quasi-periodic monodromy on $E_\tau$
(Theorem~\ref{thm:elliptic-spectral-dichotomy}).  Since
$c_0 = \partial T \ne 0$ for Virasoro, the curvature and
braiding entangle at genus~$1$: $r_1(z)$ depends on the
modular parameter $\tau$ through the Weierstrass functions.

\item \emph{Complementarity:}
$\kappa(\mathrm{Vir}_c) + \kappa(\mathrm{Vir}_{26-c}) = c/2 + (26-c)/2 = 13$
(Volume~I, Theorem~C): the Virasoro complementarity constant
is $13$, not $0$.  The complementarity potential
$F_1(\mathrm{Vir}_c) + F_1(\mathrm{Vir}_{26-c})
= \kappa/24 + \kappa'/24 = 13/24$
does not vanish: the bulk gravitational theory carries a net
genus-$1$ anomaly of $13/24$.
\end{enumerate}
\end{proposition}

\begin{proof}
(i) is Theorem~\ref{thm:virasoro-r-matrix}.
(ii) is Remark~\ref{rem:yangian-virasoro-nonformality}: the
Virasoro has infinite shadow depth $r_{\max} = \infty$ (Volume~I).
(iii) follows from Theorem~\ref{thm:elliptic-spectral-dichotomy}(c)
with $c_0 = \partial T \ne 0$.
(iv) is a direct computation from Volume~I,
$\kappa(\mathrm{Vir}_c) = c/2$ (the scalar curvature).
\end{proof}

\begin{remark}[The gravitational Yangian and $\mathrm{Vir}_{26-c}$
as resonance shadow]
% label removed: rem:gravitational-yangian-resonance
\index{resonance shadow!Virasoro}
The Koszul dual $\mathrm{Vir}_{26-c}$ is the depth-zero resonance
shadow of $\mathrm{Vir}_c$ (Volume~I,
Remark~\ref*{V1-rem:virasoro-resonance-model}): it is the image of
the finite-dimensional resonance truncation $R_{\mathrm{Vir}}$,
not the final Koszul dual.  The genuine $\mathcal{W}_\infty$-type
dual requires the full resonance-filtered completion
(Volume~I, Theorem~\ref*{V1-thm:platonic-completion}).  The
gravitational Yangian $Y_T^{\mathrm{mod}}(\mathrm{Vir}_c)$ is
therefore defined using the resonance shadow as a first
approximation; the full gravitational quantum group requires
completion technology from MC4.
\end{remark}


%%%%%%%%%%%%%%%%%%%%%%%%%%%%%%%%%%%%%%%%%%%%%%%%%%%%%%%%%%%%%%%%%%%%%%%%%%%%%
% HIGHER-GENUS BRAIDING
%%%%%%%%%%%%%%%%%%%%%%%%%%%%%%%%%%%%%%%%%%%%%%%%%%%%%%%%%%%%%%%%%%%%%%%%%%%%%

\section{Higher-genus braiding}
% label removed: sec:higher-genus-braiding
\index{higher-genus braiding|textbf}
\index{spectral braiding!higher genus}

The spectral $R$-matrix of \S\ref{subsec:YBE} is defined at
genus~$0$, where the holomorphic propagator is $dz/z$ on $\C^\times$.
At genus $g \ge 1$, the propagator acquires quasi-periodicity from
the period matrix of the curve $\Sigma_g$, forcing the braiding
into the framework of meromorphic functions on $\Sigma_g$ with
prescribed automorphy.  This section develops the genus-$1$ case
explicitly and describes the structure at higher genera.

\subsection{Genus-$1$ spectral braiding}
% label removed: subsec:genus-1-braiding

\subsubsection{The Atiyah extension and meromorphic functions on
the torus}
\index{Atiyah extension!genus 1}

Let $E_\tau = \C / (\Z + \Z\tau)$ be the elliptic curve with
modular parameter $\tau \in \mathfrak{H}$ (upper half-plane).
The holomorphic line bundle $\cO_{E_\tau}(p_0)$ (sections with at
most a simple pole at the origin $p_0$) fits into the
\emph{Atiyah extension}
\begin{equation}
% label removed: eq:atiyah-extension
0 \to \cO_{E_\tau} \to \mathcal{E}_\tau \to
T_{E_\tau} \to 0,
\end{equation}
where $\mathcal{E}_\tau$ is the Atiyah bundle and $T_{E_\tau}$ is
the tangent bundle.  The Atiyah class
$a(\cO(p_0)) \in H^1(E_\tau, \Omega^1_{E_\tau})$ is the
obstruction to splitting~\eqref{eq:atiyah-extension} and equals
the class of the Arakelov form $\omega_1$.

The spectral $r$-matrix at genus~$1$ is a section of
$\mathcal{E}_\tau \boxtimes \mathcal{E}_\tau(-\Delta)$
(meromorphic with poles along the diagonal $\Delta$ of
$E_\tau \times E_\tau$), and its quasi-periodicity is controlled
by the Atiyah extension.  The non-splitting of~\eqref{eq:atiyah-extension}
is the geometric source of the curvature-braiding entanglement
(Theorem~\ref{thm:elliptic-spectral-dichotomy}).

\subsubsection{The KZB connection as genus-$1$ shadow connection}
\index{KZB connection|textbf}
\index{shadow connection!genus 1}

\begin{definition}[Genus-$1$ shadow connection]
% label removed: def:genus-1-shadow-connection
The \emph{genus-$1$ shadow connection} is the flat connection on
the configuration space $E_\tau^n \setminus \Delta$ given by
\begin{equation}
% label removed: eq:kzb-connection
\nabla_i^{\mathrm{KZB}}
\;=\;
\partial_{z_i}
\;-\;
\sum_{j \ne i}
r^{E_\tau}_{ij}(z_i - z_j),
\end{equation}
where $r^{E_\tau}(z)$ is the elliptic $r$-matrix
of~\eqref{eq:elliptic-r-matrix}.  This is the
\emph{Knizhnik--Zamolodchikov--Bernard (KZB) connection}
\cite{Ber88}, recognized here as the genus-$1$ projection of the
shadow connection
$\nabla^{\mathrm{sh}}_{g,n} = d - \Sh_{g,n}(\Theta_\cA)$
from the universal Maurer--Cartan element
(Volume~I, Construction~\ref*{V1-const:vol1-modular-tangent-complex}).
\end{definition}

\begin{theorem}[KZB flatness from the genus-$1$ MC equation;
\ClaimStatusProvedHere]
% label removed: thm:kzb-flatness-from-mc
\index{KZB connection!flatness}
The KZB connection~\eqref{eq:kzb-connection} is flat:
$[\nabla_i^{\mathrm{KZB}}, \nabla_j^{\mathrm{KZB}}] = 0$ for
all $i \ne j$.  The flatness is equivalent to the genus-$1$
Maurer--Cartan equation for $r^{E_\tau}$ in the pronilpotent
Yangian.
\end{theorem}

\begin{proof}
\noindent\textbf{Step 1: Curvature computation.}
The curvature of $\nabla^{\mathrm{KZB}}$ on
$E_\tau^n \setminus \Delta$ is
\begin{align*}
[\nabla_i, \nabla_j]
&= -\partial_{z_j}\, r_{ij}^{E_\tau}(z_{ij})
+ \partial_{z_i}\, r_{ij}^{E_\tau}(z_{ij})
+ \sum_{k \ne i,j}
\bigl([r_{ik}^{E_\tau}(z_{ik}), r_{jk}^{E_\tau}(z_{jk})]
\bigr) \\
&\quad + [r_{ij}^{E_\tau}(z_{ij}), \sum_{k \ne i,j}
(r_{ik}^{E_\tau}(z_{ik}) + r_{jk}^{E_\tau}(z_{jk}))].
\end{align*}
The first line vanishes by the CYBE for $r^{E_\tau}$ (the
elliptic CYBE, which holds by the same Stokes argument as
Remark~\ref{rem:cybe-arnold} applied to $\FM_3(E_\tau)$).
The second line vanishes because $r^{E_\tau}$ satisfies the
\emph{classical dynamical Yang--Baxter equation} (CDYBE):
\begin{equation}
% label removed: eq:cdybe
\partial_{z_j} r_{ij}^{E_\tau}(z_{ij})
+ [r_{ij}^{E_\tau}(z_{ij}),
r_{ik}^{E_\tau}(z_{ik}) + r_{jk}^{E_\tau}(z_{jk})]
+ [r_{ik}^{E_\tau}(z_{ik}), r_{jk}^{E_\tau}(z_{jk})] = 0,
\end{equation}
which is the genus-$1$ analogue of the CYBE, incorporating the
$\tau$-dependence of $r^{E_\tau}$ via the $z$-derivative term.

\medskip
\noindent\textbf{Step 2: Stokes' theorem on $\FM_3(E_\tau)$.}
The geometric proof: apply Stokes' theorem to the closed logarithmic
form $\omega_3^{E_\tau}$ on $\FM_3(E_\tau)$, the Fulton--MacPherson
compactification of three points on the elliptic curve.  The boundary
strata of $\FM_3(E_\tau)$ are the same as for $\FM_3(\C)$: three
pairwise collision divisors $D_{\{1,2\}}$, $D_{\{1,3\}}$,
$D_{\{2,3\}}$ and the full-collision stratum $D_{\{1,2,3\}}$, but
the weight forms on each stratum now involve the elliptic propagator
$\zeta(z|\tau)$ instead of $1/z$.

The key difference from genus~$0$: the form $\omega_3^{E_\tau}$ is
NOT globally closed on $\Conf_3(E_\tau)$.  It satisfies
$d\omega_3^{E_\tau} = 2\eta_\tau \cdot c_0 \cdot \omega_1$
(the Arakelov defect from the quasi-periodicity of $\zeta$,
cf.\ Theorem~\ref{thm:elliptic-spectral-dichotomy}).  Stokes'
theorem therefore gives
\[
\int_{\FM_3(E_\tau)} d\omega_3^{E_\tau}
= \sum_{\text{faces}} \int_{\text{face}} \omega_3^{E_\tau}
= 2\eta_\tau \cdot c_0 \cdot
\int_{\FM_3(E_\tau)} \omega_1.
\]
When $c_0 = 0$ (Cartan type), the defect vanishes and the boundary
terms cancel as at genus~$0$, giving the elliptic CYBE directly.
When $c_0 \ne 0$ (Yangian type), the defect is absorbed into the
$z$-derivative term $\partial_{z_j} r_{ij}^{E_\tau}$ of the
CDYBE~\eqref{eq:cdybe}: the quasi-periodicity of $\zeta$ produces
both the monodromy defect and the dynamical correction, and these
cancel in the full CDYBE.

\medskip
\noindent\textbf{Step 3: MC equation identification.}
The CDYBE~\eqref{eq:cdybe} is the genus-$1$ MC equation
$d\, r_1 + [r_0, r_1] + \frac{1}{2}[r_0, r_0]^{g=1} = 0$
of the pronilpotent Yangian
(Theorem~\ref{thm:stable-graph-ybe-mc}, Step~3), with:
\begin{itemize}
\item $d\, r_1$ corresponding to $\partial_{z_j}
r_{ij}^{E_\tau}$ (the derivative term);
\item $[r_0, r_1]$ corresponding to
$[r_{ij}^{E_\tau}, r_{ik}^{E_\tau} + r_{jk}^{E_\tau}]$
(the mixed tree/one-loop bracket);
\item $\frac{1}{2}[r_0, r_0]^{g=1}$ corresponding to
$[r_{ik}^{E_\tau}, r_{jk}^{E_\tau}]$ (the self-loop
contribution from the genus-$0$ $r$-matrix sewn to itself
along one handle).
\end{itemize}
\end{proof}

\subsubsection{Derivation of the genus-$1$ $r$-matrix from the
Bergman kernel}
\index{Bergman kernel!genus-1 r-matrix}

\begin{theorem}[Genus-$1$ $r$-matrix; \ClaimStatusProvedHere]
% label removed: thm:genus-1-r-matrix
For the affine Kac--Moody algebra $\widehat{\fg}_k$ on the
elliptic curve $E_\tau$, the genus-$1$ classical $r$-matrix is
\begin{equation}
% label removed: eq:genus-1-r-matrix
\boxed{
r_1(z, \tau)
\;=\;
\Omega \cdot \zeta(z|\tau)
\;+\;
k\,\kappa \cdot \wp(z|\tau),
}
\end{equation}
where $\zeta(z|\tau)$ and $\wp(z|\tau)$ are the Weierstrass
zeta and $\wp$-functions with periods $1, \tau$.  This is the
elliptic $r$-matrix of~\eqref{eq:elliptic-r-matrix} specialized
to $\widehat{\fg}_k$.
\end{theorem}

\begin{proof}
By Theorem~\ref{thm:elliptic-spectral-dichotomy}, the genus-$1$
$r$-matrix is obtained from the genus-$0$ $r$-matrix
$r_0(z) = \Omega/((k+h^\vee)z) + k\kappa/z^2$ by replacing the
rational kernel functions by their elliptic counterparts:
$1/z \leadsto \zeta(z|\tau)$ and $1/z^2 \leadsto \wp(z|\tau)$.

The Bergman kernel on $E_\tau$ is
$B(z, w|\tau) = \wp(z-w|\tau)\, dz\, dw + 2\eta_1\, dz\, dw$,
where the correction $2\eta_1$ is the Legendre constant ensuring
the normalization $\int_{A} B(z, w) = 0$ (integration over the
$A$-cycle in either variable).  The one-particle restriction
operator (Volume~I, \S\ref*{sec:heisenberg-one-particle-sewing})
has eigenvalues $\wp^{(n)}(z|\tau)$ acting on the $n$-th
derivative of the current.

For the $1/z$ term (from the structure constants $f^{ab}_c$):
the unique meromorphic function on $E_\tau$ with a simple pole
at $z = 0$ and residue~$1$ is $\zeta(z|\tau)$ (up to a constant,
which is fixed by the condition that the $A$-cycle average
vanishes).  The quasi-periodicity
$\zeta(z+\tau|\tau) = \zeta(z|\tau) + 2\eta_\tau$ is the
obstruction to $r_1$ being a true meromorphic function on $E_\tau$:
the Atiyah class is nonvanishing.

For the $1/z^2$ term (from the level $k$): $\wp(z|\tau)$ is
doubly periodic with a double pole at $z = 0$, so this
contribution is a genuine meromorphic function on $E_\tau$.

Combining: $r_1(z, \tau) = \Omega \cdot \zeta(z|\tau) + k\kappa
\cdot \wp(z|\tau)$, where we absorb the normalization
$1/(k+h^\vee)$ into $\Omega$ for the simplified
expression~\eqref{eq:genus-1-r-matrix}.
\end{proof}

\begin{proposition}[Modular transformation of $r_1$;
\ClaimStatusProvedHere]
% label removed: prop:modular-transformation-r1
\index{modular transformation!r-matrix}
Under the modular $S$-transformation $\tau \mapsto -1/\tau$, the
genus-$1$ $r$-matrix transforms as
\begin{equation}
% label removed: eq:modular-S-r1
r_1(z/\tau,\, -1/\tau)
\;=\;
\tau^2\, r_1(z, \tau)
\;-\;
2\pi i\, \Omega \cdot z,
\end{equation}
where the first term is the weight-$2$ transformation of a
quasi-modular object and the second term is the anomalous
correction from the quasi-periodicity of $\zeta$.
\end{proposition}

\begin{proof}
The Weierstrass functions transform under $\tau \mapsto -1/\tau$,
$z \mapsto z/\tau$ as:
\begin{align}
\zeta(z/\tau | -1/\tau)
&= \tau\, \zeta(z|\tau) - 2\eta_\tau\, z/\tau
= \tau\, \zeta(z|\tau) + 2\pi i\, z/\tau
+ (\text{lower order in }E_2), % label removed: eq:zeta-S-transform \\
\wp(z/\tau | -1/\tau)
&= \tau^2\, \wp(z|\tau). % label removed: eq:wp-S-transform
\end{align}
The $\wp$-function transforms as a modular form of weight~$2$
(doubly periodic); the $\zeta$-function transforms anomalously,
with the quasi-period $2\eta_\tau$ producing the correction.
Using the Legendre relation $\eta_\tau \cdot 1 - \eta_1 \cdot \tau
= \pi i$ and the modular transformation of $\eta$-values:
$\eta_{-1/\tau} = \tau\,\eta_\tau$, $\eta_1' = \tau^{-1}\eta_1
- \tau^{-1}\pi i$, the full transformation gives:
\begin{align*}
r_1(z/\tau, -1/\tau)
&= \Omega \cdot \zeta(z/\tau | -1/\tau) + k\kappa \cdot
\wp(z/\tau | -1/\tau) \\
&= \Omega \cdot (\tau\,\zeta(z|\tau) + 2\pi i\, z/\tau)
+ k\kappa \cdot \tau^2\, \wp(z|\tau) \\
&= \tau\bigl(\Omega \cdot \zeta(z|\tau) + k\kappa\tau\,
\wp(z|\tau)\bigr) + 2\pi i\, \Omega \cdot z/\tau.
\end{align*}
After rescaling to account for the weight of $r_1$ as a
$(1,0)$-form in $z$: $r_1$ at genus~$1$ is a section of
$\Omega^1_{E_\tau}(\ast\, 0)$, and the $S$-transformation
$z \mapsto z/\tau$ introduces the factor $\tau$ on sections.
The final result is~\eqref{eq:modular-S-r1} after accounting
for the Jacobian $dz \mapsto dz/\tau$.
\end{proof}

\begin{remark}[Connection to the Eisenstein series]
% label removed: rem:eisenstein-r-matrix
\index{Eisenstein series!in genus-1 r-matrix}
The quasi-modular Eisenstein series $E_2(\tau) = 1 - 24
\sum_{n=1}^\infty \sigma_1(n)\, q^n$ (with $q = e^{2\pi i\tau}$)
enters the genus-$1$ $r$-matrix through the relation
$\eta_1 = \pi^2 E_2(\tau)/6$: the quasi-period of $\zeta$ is
proportional to $E_2(\tau)$.  The failure of $r_1$ to be modular
is thus controlled by $E_2$, which is the unique quasi-modular
form of weight~$2$.  The completion $\widehat{E}_2(\tau) =
E_2(\tau) - 3/(\pi\, \mathrm{Im}(\tau))$ (the Eisenstein--Maass
form) restores modularity at the cost of holomorphicity: the
``non-holomorphic completion'' of the $r$-matrix is
\[
\widehat{r}_1(z, \tau)
= \Omega \cdot \widehat{\zeta}(z|\tau) + k\kappa \cdot \wp(z|\tau),
\]
where $\widehat{\zeta}$ is the completed zeta function.  This is
a true section of $\Omega^1_{E_\tau}(\ast\, 0)$ on the modular
curve, but non-holomorphic.  The choice between holomorphicity
(the algebraic $r$-matrix) and modularity (the completed
$r$-matrix) is the genus-$1$ avatar of the holomorphic anomaly.
\end{remark}

\subsection{Genus-$2$ and beyond: the stable-graph expansion}
% label removed: subsec:genus-2-and-beyond

\subsubsection{Structure of the higher-genus $R$-matrix}
\index{higher-genus R-matrix!structure}

\begin{proposition}[Higher-genus $R$-matrix structure;
\ClaimStatusProvedHere]
% label removed: prop:higher-genus-r-structure
At genus $g \ge 2$, the classical $r$-matrix $r_g(z)$ is
determined by the MC equation~\eqref{eq:mc-at-genus-g} in the
pronilpotent Yangian.  The structure is:
\begin{enumerate}[label=\textup{(\roman*)}]
\item \emph{Propagator:} The genus-$g$ propagator is the
Szeg\H{o} kernel $S_g(z, w|\Omega_g)$ on the Riemann surface
$\Sigma_g$ of genus $g$, where $\Omega_g$ is the period matrix.
Near the diagonal $z = w$, $S_g(z,w) = 1/(z-w) + O(1)$; the
higher-order terms depend on $\Omega_g$ through theta functions.

\item \emph{Graph sum:} $r_g(z)$ is a sum over stable graphs
$\Gamma$ of genus $g$ with two external legs:
\begin{equation}
% label removed: eq:higher-genus-r-graph-sum
r_g(z)
= \sum_{\Gamma \in \StGraph_{g,2}}
\frac{1}{|\Aut(\Gamma)|}\,
\ell_\Gamma(z),
\end{equation}
where $\ell_\Gamma(z)$ is the amplitude of $\Gamma$, computed
by integrating the product of edge propagators $S_g(z_e|\Omega_g)$
over the internal vertices using the OPE data of~$\cA$.

\item \emph{Period-matrix dependence:} $r_g$ depends on the
period matrix $\Omega_g$ of $\Sigma_g$ and is a section of
a bundle on $\mathcal{A}_g$ (the moduli of principally polarized
abelian varieties) pulled back to $\mathcal{M}_g$.

\item \emph{Degeneration:} Under the clutching map
$\overline{\mathcal{M}}_{g_1, n_1+1} \times
\overline{\mathcal{M}}_{g_2, n_2+1} \to
\overline{\mathcal{M}}_{g, n}$ (with $g = g_1 + g_2$,
$n = n_1 + n_2$), the amplitude $\ell_\Gamma$ factorizes as
\[
\ell_\Gamma \;\to\;
\ell_{\Gamma_1} \cdot r_0(z_e) \cdot \ell_{\Gamma_2},
\]
where $\Gamma_1, \Gamma_2$ are the two subgraphs separated
by the edge $e$ at the node, and $r_0(z_e)$ is the genus-$0$
$r$-matrix at the spectral parameter of the node.  This
factorization is the clutching law for the $R$-matrix.
\end{enumerate}
\end{proposition}

\begin{proof}
(i) The propagator on $\Sigma_g$ is the Green's function of
the $\dbar$-operator, which is the Szeg\H{o} kernel (for spinor
fields) or the Bergman kernel (for $(1,0)$-forms).  Near the
diagonal, the local structure $1/(z-w) + O(1)$ is universal.

(ii) The MC equation~\eqref{eq:mc-at-genus-g} is a sum over
stable graphs by definition of the graph-composition bracket
(Definition~\ref{def:lie-bracket-graph-composition}).  Each graph
$\Gamma$ with $g$ loops and two external legs contributes
$\ell_\Gamma(z)/|\Aut(\Gamma)|$ to $r_g$.

(iii) The amplitudes $\ell_\Gamma$ are integrals of holomorphic
functions on $\Sigma_g$, which depend on the complex structure
through the period matrix $\Omega_g$.  More precisely, the
propagator $S_g$ is a theta function, and the integrals produce
Siegel modular forms on $\mathcal{A}_g$; pulling back to
$\mathcal{M}_g$ via the Torelli map gives the dependence on the
moduli of curves.

(iv) The clutching law follows from the factorization property
of the bar complex on degenerate curves (Volume~I,
Construction~\ref*{V1-const:vol1-graphwise-log-fm-cocomposition}):
the amplitude of a graph on a reducible curve factors through the
amplitudes on each component, with the propagator at the node.
\end{proof}

\begin{proposition}[Formal convergence criterion for the stable-graph expansion;
\ClaimStatusProvedHere]
% label removed: prop:stable-graph-convergence
\index{stable-graph expansion!convergence}
Assume a candidate series
$R_T^{\mathrm{mod}}(z; \hh) = \sum_{g \ge 0} \hh^{2g}\, r_g(z)$
has coefficients $r_g(z)$ defined by stable-graph amplitudes for a
chiral algebra $\cA$ satisfying the HS-sewing condition
\textup{(}Volume~I,
Theorem~\textup{\ref*{thm:general-hs-sewing}):} polynomial OPE
growth and subexponential sector growth\textup{)}, the genus
expansion~\eqref{eq:stable-graph-r-matrix}
converges in the $\hh$-adic topology on $Y_T^{\mathrm{mod}}$.
The convergence is absolute: the series
$\sum_g |\hh|^{2g}\, \|r_g\|$ converges for $|\hh|$
sufficiently small, where $\|r_g\|$ is the operator norm
of $r_g$ acting on finite-dimensional modules.
\end{proposition}

\begin{proof}
The HS-sewing condition (polynomial OPE growth $\|C_n\| \le
P(n)$ for a polynomial $P$, and subexponential sector growth
$\dim \cA_n \le e^{n^\alpha}$ for $\alpha < 1$) implies that
the graph amplitudes $\ell_\Gamma$ are bounded by
$|\ell_\Gamma| \le C^{|V(\Gamma)|}\, D^{|E(\Gamma)|}$ for
constants $C, D$ depending on the OPE data and the propagator
norms.  The number of stable graphs $\Gamma$ of genus $g$ with
$n$ external legs is bounded by $(Cg)^{Cg}$ for a constant $C$
(the enumeration of stable graphs by genus).  The series
$\sum_g |\hh|^{2g}\, \sum_\Gamma |\ell_\Gamma|$ converges for
$|\hh| < 1/D'$ by comparison with the exponential generating
function, where $D'$ depends on the propagator norms and the
OPE bounds.

The HS-sewing condition is verified for the entire standard
landscape by Volume~I, Theorem~\ref*{thm:general-hs-sewing}:
polynomial OPE growth holds for all Lie-theoretic chiral algebras
(the OPE coefficients grow polynomially in the mode number),
and subexponential sector growth holds for all chiral algebras
with polynomial character asymptotics (Heisenberg, affine, Virasoro,
$\mathcal{W}$-algebras at generic parameters).
\end{proof}


%%%%%%%%%%%%%%%%%%%%%%%%%%%%%%%%%%%%%%%%%%%%%%%%%%%%%%%%%%%%%%%%%%%%%%%%%%%%%
% HOPF STRUCTURE AND REPRESENTATION THEORY
%%%%%%%%%%%%%%%%%%%%%%%%%%%%%%%%%%%%%%%%%%%%%%%%%%%%%%%%%%%%%%%%%%%%%%%%%%%%%

\section{Hopf structure and representation theory}
% label removed: sec:hopf-structure
\index{Hopf structure!pronilpotent Yangian|textbf}
\index{representation theory!pronilpotent Yangian}

The pronilpotent Yangian $Y_T^{\mathrm{mod}}$ carries a Hopf-like
structure inherited from the sewing/clutching maps on
$\overline{\mathcal{M}}_{g,n}$.  This section constructs the
coproduct, antipode, and representation theory, connecting the
abstract MC framework to concrete physical observables.

\subsection{The coproduct from the separating clutching map}
% label removed: subsec:coproduct-clutching

\begin{definition}[Coproduct via separating degeneration]
% label removed: def:coproduct-separating
\index{coproduct!separating degeneration}
\index{clutching map!separating}
The \emph{separating clutching map} is the morphism
\begin{equation}
% label removed: eq:separating-clutching
\xi_{\mathrm{sep}}^{g_1, g_2} \colon
\overline{\mathcal{M}}_{g_1, n_1+1}
\times
\overline{\mathcal{M}}_{g_2, n_2+1}
\longrightarrow
\overline{\mathcal{M}}_{g, n},
\end{equation}
with $g = g_1 + g_2$ and $n = n_1 + n_2$, defined by identifying
the last marked point of the first curve with the last marked point
of the second curve to form a nodal curve of genus $g$ with $n$
marked points.  The separating node separates the two components.

The \emph{coproduct} on $Y_T^{\mathrm{mod}}$ is the linear map
\begin{equation}
% label removed: eq:coproduct
\Delta \colon Y_T^{\mathrm{mod}}
\longrightarrow
Y_T^{\mathrm{mod}} \;\widehat{\otimes}\; Y_T^{\mathrm{mod}},
\end{equation}
defined by pullback along $\xi_{\mathrm{sep}}$: for
$\alpha = \{\alpha_{g,n}\} \in Y_T^{\mathrm{mod}}$,
\[
\Delta(\alpha)_{g_1, n_1;\, g_2, n_2}
= (\xi_{\mathrm{sep}}^{g_1, g_2})^\ast(\alpha_{g_1+g_2, n_1+n_2}).
\]
The completed tensor product $\widehat{\otimes}$ is with respect
to the arity filtration on each factor.
\end{definition}

\begin{theorem}[Coassociativity from composition of clutching maps;
\ClaimStatusProvedHere]
% label removed: thm:coassociativity
\index{coassociativity!from clutching composition}
The coproduct~\eqref{eq:coproduct} is coassociative:
\begin{equation}
% label removed: eq:coassociativity
(\Delta \otimes \id) \circ \Delta
\;=\;
(\id \otimes \Delta) \circ \Delta
\colon
Y_T^{\mathrm{mod}}
\to
Y_T^{\mathrm{mod}}
\;\widehat{\otimes}\;
Y_T^{\mathrm{mod}}
\;\widehat{\otimes}\;
Y_T^{\mathrm{mod}}.
\end{equation}
\end{theorem}

\begin{proof}
Both sides of~\eqref{eq:coassociativity} are pullbacks of
$\alpha_{g, n}$ along the iterated clutching map
$\overline{\mathcal{M}}_{g_1, n_1+1} \times
\overline{\mathcal{M}}_{g_2, n_2+1} \times
\overline{\mathcal{M}}_{g_3, n_3+1}
\to \overline{\mathcal{M}}_{g, n}$
(with $g = g_1 + g_2 + g_3$, $n = n_1 + n_2 + n_3$), which
identifies the last point of the first curve with the last point
of the second, and the last point of the second with the last
point of the third.

The two sides correspond to the two orderings of the
identification:
\begin{itemize}
\item $(\Delta \otimes \id) \circ \Delta$: first identify
$\Sigma_{g_1} \cup \Sigma_{g_2+g_3}$, then split
$\Sigma_{g_2+g_3}$ into $\Sigma_{g_2} \cup \Sigma_{g_3}$.
\item $(\id \otimes \Delta) \circ \Delta$: first identify
$\Sigma_{g_1+g_2} \cup \Sigma_{g_3}$, then split
$\Sigma_{g_1+g_2}$ into $\Sigma_{g_1} \cup \Sigma_{g_2}$.
\end{itemize}
Both produce the same 3-component nodal curve
$\Sigma_{g_1} \cup \Sigma_{g_2} \cup \Sigma_{g_3}$ with two
separating nodes, so the pullbacks agree.  The verification
at the chain level uses the homotopy commutativity of the
clutching maps: the two orderings define homotopic maps
$\overline{\mathcal{M}}_{g_1} \times \overline{\mathcal{M}}_{g_2}
\times \overline{\mathcal{M}}_{g_3} \to
\overline{\mathcal{M}}_g$ because the order of node formation
is irrelevant in the stable category (the universal curve is a
functor from the nerve of the combinatorial category of stable
graphs, and the nerve is a Kan complex).
\end{proof}

\subsection{The antipode from the involution on stable graphs}
% label removed: subsec:antipode

\begin{definition}[Antipode via duality]
% label removed: def:antipode
\index{antipode!from stable graph involution}
The \emph{antipode} on $Y_T^{\mathrm{mod}}$ is the linear map
\begin{equation}
% label removed: eq:antipode
S \colon Y_T^{\mathrm{mod}}
\longrightarrow Y_T^{\mathrm{mod}}
\end{equation}
defined by the involution on stable graphs that reverses the
orientation of all edges.  Concretely, for
$\alpha_{g,n} \in \Hom(C_\ast(\overline{\mathcal{M}}_{g,n+1})
\otimes (\cA^!)^{\otimes n}, \cA^!)$:
\[
S(\alpha)_{g,n}(\gamma; a_1, \ldots, a_n)
= (-1)^{g + n + |\alpha|}\,
\alpha_{g,n}(\gamma^\vee; a_n, \ldots, a_1),
\]
where $\gamma^\vee$ is the Poincar\'e dual class in
$C^\ast(\overline{\mathcal{M}}_{g,n+1})$ and the inputs are
reversed.  The sign $(-1)^{g+n+|\alpha|}$ is determined by the
Koszul convention on the desuspended bar complex.
\end{definition}

\begin{proposition}[Antipode axioms; \ClaimStatusProvedHere]
% label removed: prop:antipode-axioms
\index{antipode!axioms}
The antipode satisfies the standard Hopf-algebra axioms (in the
completed, pronilpotent sense):
\begin{enumerate}[label=\textup{(\roman*)}]
\item $m \circ (\id \otimes S) \circ \Delta = \eta \circ
\varepsilon$, where $m$ is the multiplication (graph composition),
$\eta$ is the unit, and $\varepsilon$ is the counit.
\item $m \circ (S \otimes \id) \circ \Delta = \eta \circ
\varepsilon$.
\item $S^2 = \id$ (involutivity, from the involutivity of
edge reversal on stable graphs).
\end{enumerate}
\end{proposition}

\begin{proof}
(i)-(ii): The composition $m \circ (\id \otimes S) \circ \Delta$
applied to $\alpha_{g,n}$ sums over all ways to split the genus-$g$
curve into two components at a separating node, apply $\alpha$ on
one component and $S(\alpha)$ on the other, and recompose via the
multiplication $m$ (edge contraction).  The edge-reversal involution
$S$ combined with the splitting $\Delta$ and the contraction $m$
produces a telescoping sum: at each separating node, the two
contributions from $\alpha$ and $S(\alpha)$ cancel except at the
trivial splitting $g_1 = 0, n_1 = 0$ (where one component is a
point), leaving only the counit.

(iii): Reversing all edges twice returns to the original orientation.
The sign $(-1)^{g+n+|\alpha|}$ applied twice gives $(-1)^{2(g+n+|\alpha|)}
= 1$.
\end{proof}

\subsection{Representations: highest-weight modules and line operators}
% label removed: subsec:representations

\begin{definition}[Module over the pronilpotent Yangian]
% label removed: def:yangian-module
\index{Yangian module!pronilpotent}
A \emph{module} over $Y_T^{\mathrm{mod}}$ is a collection of linear
maps
\[
\rho_{g,n} \colon
C_\ast(\overline{\mathcal{M}}_{g,n+1})
\otimes (\cA^!)^{\otimes n}
\longrightarrow \End(V)
\]
for a graded vector space $V$, compatible with the graph-composition
bracket: $\rho([\alpha, \beta]) = [\rho(\alpha), \rho(\beta)]$.
A module is \emph{highest-weight} if $V$ has a cyclic vector
$v_0$ (the highest-weight vector) annihilated by all ``positive''
convolution elements (those in arity $n \ge 2$ at genus~$0$).
\end{definition}

\begin{theorem}[Line operators on the genus-zero Yangian lane;
\ClaimStatusProvedHere]
% label removed: thm:line-operators-yangian-modules
\index{line operators!as Yangian modules}
Let\/ $\cA$ be a logarithmic\/ $\SCchtop$-algebra.  The category $\mathcal{C}_{\mathrm{line}}$
of boundary line operators (Theorem~\ref{thm:lines_as_modules}) is
equivalent to the category of modules over $Y_T^{\mathrm{mod}}$
supported at genus~$0$:
\begin{equation}
% label removed: eq:line-ops-yangian-mods
\mathcal{C}_{\mathrm{line}}
\;\simeq\;
Y_T^{\mathrm{mod}}|_{g=0}\text{-}\mathbf{mod}
\;\simeq\;
\cA^!\text{-}\mathbf{mod}.
\end{equation}
The genus-$0$ truncation of $Y_T^{\mathrm{mod}}$ acts on line
operators via the Yangian product (open-color operations), and the
spectral $R$-matrix $R(z)$ provides the braiding on this genus-zero
lane.  Any higher-genus spectral correction belongs to the separate
conjectural modular-kernel package and is not part of this theorem.
\end{theorem}

\begin{proof}
The equivalence
$\mathcal{C}_{\mathrm{line}} \simeq \cA^!_{\mathrm{line}}\text{-}\mathbf{mod}$
is Theorem~\ref{thm:lines_as_modules}.  The identification of
$\cA^!_{\mathrm{line}}$-modules with genus-$0$ modules over
$Y_T^{\mathrm{mod}}$
follows from Corollary~\ref{cor:classical-yangian-truncation}: the
genus-$0$ truncation of $Y_T^{\mathrm{mod}}$ is the classical
Yangian, and on the affine evaluation locus its representations are
representations of the open-colour dual
$\cA^!_{\mathrm{line}}$ via the standard affine identification
(Theorem~\ref{thm:Koszul_dual_Yangian}).
\end{proof}

\begin{proposition}[The evaluation representation;
\ClaimStatusProvedHere]
% label removed: prop:evaluation-representation
\index{evaluation representation|textbf}
\index{line operators!evaluation at a point}
The \emph{evaluation representation} at $z = z_0 \in \C^\times$ is
the Yangian module $\mathrm{ev}_{z_0}(V)$ defined by
\begin{equation}
% label removed: eq:evaluation-rep
\mathrm{ev}_{z_0} \colon Y_T^{\mathrm{mod}} \longrightarrow
\End(V), \qquad
\mathrm{ev}_{z_0}(\alpha)(v) = \alpha(z_0)(v),
\end{equation}
where $\alpha(z_0)$ is the evaluation of the spectral generating
function at $z = z_0$.  For a highest-weight module $V$ of
$Y(\fg)$ with highest weight $\mu$, the evaluation representation
$\mathrm{ev}_{z_0}(V)$ corresponds to the insertion of a line
operator of type $\mu$ at holomorphic position $z_0$ along the
boundary.

The evaluation intertwines the spectral braiding with the
$R$-matrix at the evaluation point:
\[
R(z_1 - z_2) \cdot (\mathrm{ev}_{z_1} \otimes \mathrm{ev}_{z_2})
= (\mathrm{ev}_{z_2} \otimes \mathrm{ev}_{z_1}) \cdot R(z_1 - z_2).
\]
This is the spectral-parameter formulation of the braiding
axiom: the $R$-matrix depends only on the difference of
evaluation points.
\end{proposition}

\begin{proof}
The evaluation map $\mathrm{ev}_{z_0}$ is the composite of the
genus-$0$ projection $Y_T^{\mathrm{mod}} \to Y(\fg)$ with the
standard evaluation $Y(\fg) \to U(\fg)$ at $z_0$ (the
specialization of the spectral parameter).  It is an algebra
homomorphism because the Yangian product is continuous in
the spectral variable, and specialization is multiplicative.

The intertwining property follows from the translation invariance
of $R(z)$: the braiding depends on $z = z_1 - z_2$, so evaluating
both factors at $z_1, z_2$ and braiding gives the same result as
braiding the formal generating functions and then evaluating.
\end{proof}

\begin{remark}[Virasoro: conical defects and BTZ black holes]
% label removed: rem:virasoro-conical-btz
\index{Virasoro algebra!conical defects}
\index{BTZ black hole!as Yangian module}
For the gravitational Yangian
$Y_T^{\mathrm{mod}}(\mathrm{Vir}_c)$, the highest-weight modules
have a gravitational interpretation:
\begin{itemize}
\item \emph{Light representations} ($h < c/24$): conical defects
in the $\mathrm{AdS}_3$ bulk.  The line operator creates a
conical singularity with deficit angle $2\pi(1 - \sqrt{1 - 24h/c})$.
The evaluation at $z_0$ inserts the defect at holomorphic position
$z_0$ along the boundary.

\item \emph{Heavy representations} ($h > c/24$): BTZ black holes.
The line operator creates a BTZ geometry with horizon area
$4\pi\sqrt{24h/c - 1}$ (in units of the $\mathrm{AdS}_3$ radius).
The spectral braiding of two BTZ modules encodes the gravitational
scattering amplitude at tree level.

\item \emph{Threshold} ($h = c/24$): the massless BTZ black hole,
the boundary between conical defects and black holes.  At this
threshold, the spectral $r$-matrix has a double pole (from the
Virasoro quartic singularity), and the braiding develops a
logarithmic singularity.
\end{itemize}
These representations span the module category of the gravitational
Yangian at genus~$0$.  At genus~$1$, the modular parameter $\tau$
enters through the KZB connection
(Definition~\ref{def:genus-1-shadow-connection}), and the BTZ
black hole becomes a thermal state at temperature
$T = 1/(2\pi\, \mathrm{Im}(\tau))$.
\end{remark}

\subsection{Module categories and the MC3 connection}
% label removed: subsec:mc3-connection
\index{MC3!module category connection}

\begin{theorem}[$\mathcal{C}_{\mathrm{line}} \simeq
\cA^!\text{-}\mathbf{mod}$: unconditional;
\ClaimStatusProvedHere]
% label removed: thm:cline-amod-unconditional
\index{line operators!as Koszul dual modules}
Let\/ $\cA$ be a logarithmic\/ $\SCchtop$-algebra and homotopy-Koszulity of $\SCchtop$
\textup{(}Theorem~\textup{\ref{thm:homotopy-Koszul})}.
The category of boundary line operators is equivalent to the
category of modules over the Koszul dual:
\begin{equation}
% label removed: eq:cline-amod
\mathcal{C}_{\mathrm{line}} \;\simeq\; \cA^!\text{-}\mathbf{mod}.
\end{equation}
This is unconditional (given the standing hypotheses) and follows
from the two-color Koszul duality
\textup{(}Theorem~\textup{\ref{thm:two-color-master}(ii))}.
\end{theorem}

\begin{proof}
The open-color Quillen equivalence of
Theorem~\ref{thm:two-color-master}(ii) identifies line operators
(right $\SCchtop$-modules supported on $\C \times \{0\}$) with
$\cA^!$-modules via the bar-cobar adjunction in the open color.
The homotopy-Koszulity ensures this is a genuine equivalence (not
just an adjunction): the counit and unit of the adjunction are
quasi-isomorphisms by the homotopy-Koszulity criterion
(Theorem~\ref{thm:homotopy-Koszul}, which is proved).
\end{proof}

\begin{remark}[Connection to MC3: all simple types]
% label removed: rem:mc3-type-a-connection
\index{MC3!all simple types connection}
For $\cA = \widehat{\fg}_k$ with $\fg$ simple, the module category
$\cA^!\text{-}\mathbf{mod}$ is the category of modules over the
Yangian $Y(\fg)$.  On the current theorem surface, the all-types MC3
results provide the categorical Clebsch--Gordan package and the DK
comparison on the evaluation-generated core.  In type~$A$, the later
shifted-prefundamental and pro-Weyl packets further reduce the
post-core gap to the compact-completion conjecture; in general type
the analogous post-core mechanism remains conditional.  Thus this
remark supplies the all-types evaluation-core connection to MC3, not a
theorematic identification of the full Yangian module category with
the completed line-side package.
\end{remark}


%%%%%%%%%%%%%%%%%%%%%%%%%%%%%%%%%%%%%%%%%%%%%%%%%%%%%%%%%%%%%%%%%%%%%%%%%%%%%
% FREDHOLM PARTITION FUNCTIONS AND SEWING
%%%%%%%%%%%%%%%%%%%%%%%%%%%%%%%%%%%%%%%%%%%%%%%%%%%%%%%%%%%%%%%%%%%%%%%%%%%%%

\section{Fredholm partition functions and sewing}
% label removed: sec:fredholm-partition-functions
\index{Fredholm determinant!partition function|textbf}
\index{sewing!Fredholm partition function}
\index{partition function!Fredholm}

The abstract machinery of the pronilpotent Yangian and the
stable-graph $R$-matrix produces explicit numbers when evaluated on
the simplest chiral algebra: the Heisenberg $\cH_k$.  This section
derives the genus-$g$ partition function of the Heisenberg algebra
from the sewing axioms, recovering the Dedekind eta function as a
Fredholm determinant.

\subsection{The sewing envelope for the Heisenberg algebra}
% label removed: subsec:sewing-envelope-heisenberg
\index{sewing envelope!Heisenberg|textbf}

\begin{definition}[Sewing envelope]
% label removed: def:sewing-envelope
\index{sewing envelope!definition}
The \emph{sewing envelope} $\cA^{\mathrm{sew}}$ of a chiral algebra
$\cA$ is the Hausdorff completion of $\cA$ with respect to the
\emph{all-amplitude seminorm topology}: the coarsest topology
making all sewing amplitudes $\langle a_1 \otimes \cdots \otimes a_n
\rangle_{g,n}$ continuous.  Elements of $\cA^{\mathrm{sew}}$ are
formal sums $\sum_\alpha c_\alpha\, a_\alpha$ for which every
sewing amplitude converges absolutely.
\end{definition}

\begin{theorem}[Heisenberg sewing envelope; \ClaimStatusProvedHere]
% label removed: thm:heisenberg-sewing-envelope
\index{sewing envelope!Heisenberg}
For the Heisenberg algebra $\cH_k$, the sewing envelope is
\begin{equation}
% label removed: eq:heisenberg-sewing-envelope
\cH_k^{\mathrm{sew}}
\;\simeq\;
\Sym\bigl(\cA^2(D)\bigr),
\end{equation}
where $\cA^2(D)$ is the space of square-summable holomorphic
functions on the unit disk $D = \{|z| < 1\}$, and $\Sym$
denotes the completed symmetric algebra \textup{(}bosonic
Fock space\textup{)}.

The sewing amplitudes at genus $g$ with $n$ insertions are
trace-class:
\[
\langle a_1 \otimes \cdots \otimes a_n \rangle_{g,n}
= \Tr_{\cH_k^{\mathrm{sew}}}
\bigl(a_1(z_1) \cdots a_n(z_n) \cdot q^{L_0 - c/24}\bigr)^g
\]
with $q = e^{2\pi i\tau}$, and the trace converges absolutely
for $|q| < 1$ by the Hilbert--Schmidt condition.
\end{theorem}

\begin{proof}
The Heisenberg algebra $\cH_k$ is generated by modes $J_n$
($n \in \Z$) with $[J_m, J_n] = k\,m\,\delta_{m+n,0}$.  The
Fock space $\mathcal{F}_k = \Sym(\bigoplus_{n \ge 1} \C \cdot J_{-n})$
is the vacuum module.  The sewing amplitude at genus~$1$ is the
partition function
$Z_1(\cH_k, \tau) = \Tr_{\mathcal{F}_k}(q^{L_0 - k/24})$,
which converges for $|q| < 1$ because $L_0$ has spectrum
$\{n_1 + n_2 + \cdots : 1 \le n_1 \le n_2 \le \cdots\}$
with finite multiplicities (the partitions of~$n$).

The Hausdorff completion of $\mathcal{F}_k$ in the sewing-amplitude
topology is $\Sym(\cA^2(D))$: the space of square-summable
symmetric tensors on the disk.  This follows from the identification
of the Heisenberg Fock space with the symmetric Fock space of the
Hardy space $H^2(D)$ via the map $J_{-n} \mapsto z^n/\sqrt{kn}$
(the creation operator becomes multiplication by $z^n$, normalized
by the commutation relation).  The Hardy space $H^2(D)$ is the
square-summable holomorphic functions on the disk, so
$\Sym(H^2(D)) = \mathcal{F}_k$ as Hilbert spaces.  The completion
$\cA^2(D) \supset H^2(D)$ includes the boundary values, giving
the sewing envelope.
\end{proof}

\subsection{The one-particle Bergman restriction operator}
% label removed: subsec:bergman-restriction
\index{Bergman kernel!restriction operator|textbf}

\begin{definition}[One-particle restriction operator]
% label removed: def:one-particle-restriction
\index{one-particle restriction operator}
Let $\Sigma_g$ be a genus-$g$ Riemann surface with local
coordinate $z$ near a marked point $p$.  The \emph{one-particle
Bergman restriction operator} is the operator
\begin{equation}
% label removed: eq:restriction-operator
R_g \colon H^2(D) \longrightarrow H^2(D)
\end{equation}
defined by the Bergman kernel on $\Sigma_g$: for
$f \in H^2(D)$,
\[
(R_g\, f)(z) = \oint_\gamma B_g(z, w)\, f(w)\, dw,
\]
where $B_g(z, w) = \partial_z \partial_w \log E(z,w) + \omega_g(z,w)$
is the Bergman kernel on $\Sigma_g$ (with $E(z,w)$ the prime
form and $\omega_g$ the canonical bidifferential), and $\gamma$ is
a contour surrounding $p$ in the local coordinate disk.
\end{definition}

\begin{proposition}[Eigenvalues and trace class;
\ClaimStatusProvedHere]
% label removed: prop:restriction-eigenvalues
\index{one-particle restriction operator!eigenvalues}
For $\Sigma_g = E_\tau$ \textup{(}genus $1$\textup{)} with the
sewing parameter $q = e^{2\pi i\tau}$:
\begin{enumerate}[label=\textup{(\roman*)}]
\item The eigenvalues of $R_1$ are $\{q^n : n \ge 1\}$, each
with multiplicity~$1$.
\item $R_1$ is trace class: $\Tr(R_1) = \sum_{n=1}^\infty q^n
= q/(1-q)$.
\item $R_1$ is Hilbert--Schmidt: $\|R_1\|_{\mathrm{HS}}^2
= \sum_{n=1}^\infty |q|^{2n} = |q|^2/(1-|q|^2) < \infty$ for
$|q| < 1$.
\end{enumerate}
\end{proposition}

\begin{proof}
(i) On $E_\tau$, the Bergman kernel in the local coordinate
near $z = 0$ has the expansion
$B_\tau(z, w) = \wp(z-w|\tau)\, dz\, dw + 2\eta_1\, dz\, dw$.
The restriction operator $R_1$ on $H^2(D)$ with basis
$\{z^{n-1}\}_{n \ge 1}$ has matrix elements
$(R_1)_{mn} = \oint \oint w^{m-1}\, B_\tau(z,w)\, z^{-(n+1)}\,
dz\, dw / (2\pi i)^2$.  Using the Fourier expansion of $\wp$
on the torus: $\wp(z|\tau) = (2\pi i)^2\bigl(-1/12 + \sum_{n=1}^\infty
n\,(q^n z + q^n z^{-1})/(1-q^n)\bigr)$ (in multiplicative
coordinates $z \to e^{2\pi iz}$), the matrix elements evaluate
to $(R_1)_{mn} = q^n\, \delta_{mn}$.  The operator is diagonal
with eigenvalues $q^n$.

(ii) $\Tr(R_1) = \sum_{n=1}^\infty q^n = q/(1-q)$, convergent
for $|q| < 1$.

(iii) $\|R_1\|_{\mathrm{HS}}^2 = \sum_{n=1}^\infty |q^n|^2
= |q|^2/(1 - |q|^2)$, finite for $|q| < 1$.
\end{proof}

\subsection{Derivation of the partition function}
% label removed: subsec:partition-function-derivation
\index{partition function!Heisenberg derivation}

\begin{theorem}[Heisenberg partition function as Fredholm
determinant; \ClaimStatusProvedHere]
% label removed: thm:heisenberg-partition-fredholm
\index{Fredholm determinant!Heisenberg partition function}
The genus-$g$ partition function of the Heisenberg algebra
$\cH_k$ on a Riemann surface $\Sigma_g$ with period matrix
$\Omega_g$ is
\begin{equation}
% label removed: eq:heisenberg-partition
\boxed{
Z_g(\cH_k)
\;=\;
\det(1 - T_q)^{-k},
}
\end{equation}
where $T_q = R_g$ is the one-particle restriction operator
\textup{(}Definition~\textup{\ref{def:one-particle-restriction})}
with sewing parameter $q$, and $\det$ denotes the Fredholm
determinant.
\end{theorem}

\begin{proof}
\noindent\textbf{Step 1: Bosonic Fock space sewing.}
The sewing axiom for the Heisenberg algebra states that the
genus-$g$ partition function is computed by tracing the sewing
operator over the Fock space:
\[
Z_g(\cH_k)
= \Tr_{\mathcal{F}_k}(\mathcal{S}_g),
\]
where $\mathcal{S}_g$ is the sewing operator corresponding
to the identification of puncture neighborhoods on $\Sigma_g$.
For a genus-$g$ surface obtained by sewing $g$ handles onto a
sphere, $\mathcal{S}_g = (R_g)^{\otimes g}$ restricted to the
diagonal (identifying the two sides of each handle).

\medskip
\noindent\textbf{Step 2: One-particle reduction.}
The Fock space $\mathcal{F}_k = \Sym^k(H^2(D))$ factors as
the $k$-fold symmetric tensor product of the one-particle
space $H^2(D)$ (the level $k$ counts the number of independent
bosons).  By the Fock space trace formula, the trace of a
second-quantized operator $\Gamma(A) = \bigoplus_{n \ge 0}
A^{\otimes n}|_{\Sym^n}$ on $\Sym(V)$ is
\[
\Tr_{\Sym(V)}(\Gamma(A)) = \det(1 - A)^{-1},
\]
valid when $A$ is trace class with $\|A\| < 1$ (the Fredholm
determinant formula for bosonic traces).  Applied to
$A = R_g$ and $V = H^2(D)$:
\[
\Tr_{\mathcal{F}_1}(\mathcal{S}_g) = \det(1 - R_g)^{-1}.
\]
For level $k$ (i.e., $k$ independent bosons), the Fock space
is $\Sym^{\otimes k}(H^2(D))$, and the trace is
$\det(1 - R_g)^{-k}$.

\medskip
\noindent\textbf{Step 3: Identification with Fredholm determinant.}
The operator $T_q = R_g$ is trace class
(Proposition~\ref{prop:restriction-eigenvalues}(ii)) with
eigenvalues $\{q^m : m \ge 1\}$ and $\|T_q\| = |q| < 1$ for
$|q| < 1$, so $\Tr(T_q^n) = \sum_{m=1}^\infty q^{mn}$.  Hence
\begin{align*}
\log\det(1 - T_q)
&= -\sum_{n=1}^\infty \frac{1}{n}\, \Tr(T_q^n)
= -\sum_{n=1}^\infty \frac{1}{n}\, \sum_{m=1}^\infty q^{mn} \\
&= -\sum_{m=1}^\infty \sum_{n=1}^\infty \frac{q^{mn}}{n}
= \sum_{m=1}^\infty \log(1 - q^m),
\end{align*}
so $\det(1 - T_q) = \prod_{m=1}^\infty (1 - q^m)$.

\medskip
\noindent\textbf{Step 4: Product formula.}
The Fredholm determinant factors as the infinite product
\[
\det(1 - T_q) = \prod_{n=1}^\infty (1 - q^n),
\]
convergent absolutely for $|q| < 1$.  Therefore
\begin{equation}
% label removed: eq:heisenberg-partition-product
Z_g(\cH_k) = \det(1 - T_q)^{-k}
= \prod_{n=1}^\infty (1 - q^n)^{-k}.
\end{equation}
\end{proof}

\subsection{Dedekind eta from the product formula}
% label removed: subsec:dedekind-eta
\index{Dedekind eta function!from Fredholm determinant|textbf}

\begin{corollary}[Dedekind eta from sewing;
\ClaimStatusProvedHere]
% label removed: cor:dedekind-eta-from-sewing
At genus $1$ with $q = e^{2\pi i\tau}$, the Heisenberg partition
function is
\begin{equation}
% label removed: eq:heisenberg-eta
Z_1(\cH_k, \tau)
= q^{-k/24}\, \prod_{n=1}^\infty (1 - q^n)^{-k}
= \eta(\tau)^{-k},
\end{equation}
where $\eta(\tau) = q^{1/24} \prod_{n=1}^\infty (1 - q^n)$ is the
Dedekind eta function.  The prefactor $q^{-k/24}$ arises from the
vacuum energy: the zero-point energy of $k$ free bosons on the
torus is $E_0 = -k/24$ \textup{(}the Casimir
energy\textup{)}, contributing $q^{E_0} = q^{-k/24}$ to the
trace.
\end{corollary}

\begin{proof}
The genus-$1$ partition function is the torus trace
$Z_1 = \Tr_{\mathcal{F}_k}(q^{L_0 - k/24})$.  The Virasoro
generator $L_0$ on the Heisenberg Fock space has spectrum
$\{n_1 + \cdots + n_r : 1 \le n_1 \le \cdots \le n_r,\,
r \ge 0\}$ (partitions), so
$\Tr(q^{L_0}) = \prod_{n=1}^\infty (1-q^n)^{-1}$ for each
boson.  The vacuum subtraction $-k/24$ gives the prefactor
$q^{-k/24}$, and the $k$ independent bosons give the $k$-th
power.  The Dedekind eta $\eta(\tau) = q^{1/24}\prod(1-q^n)$
satisfies $\eta^{-1} = q^{-1/24}\prod(1-q^n)^{-1}$, so
$Z_1 = \eta(\tau)^{-k}$.
\end{proof}

\begin{remark}[The Fredholm determinant and the modular
characteristic]
% label removed: rem:fredholm-modular-char
\index{modular characteristic!Fredholm determinant}
The genus-$1$ free energy is
$F_1(\cH_k) = -\log Z_1 = k\log\eta(\tau) = k(2\pi i\tau/24 +
\sum_{n=1}^\infty \log(1-q^n))$.  The leading term
$k\cdot 2\pi i\tau/24$ gives the modular characteristic
$\kappa(\cH_k) = k$: the coefficient of the Arakelov form
$\omega_1 = \mathrm{Im}(\tau)^{-1}\, dz \wedge d\bar{z}/(2i)$
in the curvature $\dfib^{\,2} = \kappa \cdot \omega_1$ is
exactly $k$, matching Volume~I.  The infinite product
$\prod(1-q^n)^{-k}$ is the non-perturbative completion of
the modular characteristic: it encodes all instanton corrections
to the genus-$1$ curvature.
\end{remark}

\subsection{Non-Fredholm corrections: higher-order interactions}
% label removed: subsec:non-fredholm-corrections
\index{Feynman integral!non-Fredholm corrections}

\begin{remark}[Cubic and quartic corrections for non-abelian algebras]
% label removed: rem:cubic-quartic-corrections
\index{cubic shadow!Feynman integral}
\index{quartic resonance class!Feynman integral}
For non-abelian chiral algebras, the sewing amplitudes are NOT
purely one-particle: the higher OPE interactions ($m_3$, $m_4$,
\ldots) contribute Feynman diagrams with internal vertices of
valence $\ge 3$.  These are the \emph{non-Fredholm corrections}
to the partition function.

\begin{enumerate}[label=\textup{(\roman*)}]
\item \emph{Cubic correction (affine):}
For $\widehat{\fg}_k$ at genus $g$, the cubic vertex $m_3$
(from the cubic shadow $C_{\widehat{\fg}_k}$ of Volume~I)
contributes a three-vertex Feynman integral
\[
\delta Z_g^{(3)}
= \int_{\Sigma_g^3}
C(z_1, z_2, z_3)\,
B_g(z_1, z_2)\, B_g(z_2, z_3)\, B_g(z_3, z_1)\,
\prod_i dz_i\, d\bar{z}_i,
\]
where $C$ is the cubic OPE coefficient and $B_g$ is the
Bergman kernel.  This integral is a Siegel modular form
of weight $(3, 3)$ on $\mathcal{A}_g$.  The shadow depth
$r_{\max}(\widehat{\fg}_k) = 3$ (Lie class) means no
quartic or higher corrections arise: the Feynman expansion
terminates at three vertices.

\item \emph{Quartic correction (Virasoro):}
For $\mathrm{Vir}_c$ at genus $g$, the quartic resonance
class $Q^{\mathrm{contact}}_{\mathrm{Vir}} = 10/(c(5c+22))$
(Volume~I) contributes a four-vertex Feynman integral
\[
\delta Z_g^{(4)}
= \frac{10}{c(5c+22)}\,
\int_{\Sigma_g^4}
B_g(z_1, z_2)\, B_g(z_2, z_3)\, B_g(z_3, z_4)\,
B_g(z_4, z_1)\,
\prod_i dz_i\, d\bar{z}_i.
\]
This integral is the genus-$g$ avatar of the quartic contact
invariant and is nonzero for $c \ne 0$ and $c \ne -22/5$.
The shadow depth $r_{\max}(\mathrm{Vir}_c) = \infty$ (mixed
class) means the Feynman expansion does NOT terminate: the
quintic, sextic, \ldots corrections are all nonzero, producing
an infinite perturbative series.

\item \emph{The HS-sewing criterion:}
Despite the infinite perturbative expansion, the partition
function converges for all chiral algebras in the standard
landscape: Volume~I, Theorem~\ref*{thm:general-hs-sewing}
guarantees that polynomial OPE growth and subexponential
sector growth imply absolute convergence of the sewing
amplitudes at all genera.  The Fredholm determinant $\det(1-T_q)^{-k}$
is the leading (one-particle) term; the non-Fredholm corrections
are exponentially suppressed in $|q|$.
\end{enumerate}
\end{remark}

\begin{remark}[The Heisenberg sewing theorem and one-particle
Bergman reduction]
% label removed: rem:heisenberg-sewing-theorem
\index{Heisenberg algebra!sewing theorem}
The derivation of $Z_g(\cH_k) = \det(1-T_q)^{-k}$ is a special
case of the Heisenberg sewing theorem
(Volume~I, Theorem~\ref*{thm:heisenberg-sewing}): the entire
genus-$g$ partition function reduces to a one-particle computation
because the Heisenberg has shadow depth $r_{\max} = 2$ (Gaussian
class).  The one-particle Bergman restriction operator $R_g$ captures
all the information: no multi-particle correlations contribute.

This is the content of Volume~I,
Theorem~\ref*{V1-thm:heisenberg-one-particle-sewing}: for the
Heisenberg algebra, the factorization of the partition function
through the one-particle space is \emph{exact}, not an
approximation.  The Fredholm determinant formula is the closed-form
expression for the sewing of Gaussian chiral algebras.
\end{remark}

\begin{remark}[Moriwaki's theorem and IndHilb factorization]
% label removed: rem:moriwaki-theorem
\index{Moriwaki!theorem}
\index{factorization homology!IndHilb-valued}
The sewing envelope $\cH_k^{\mathrm{sew}} \simeq \Sym(\cA^2(D))$
of Theorem~\ref{thm:heisenberg-sewing-envelope} is a special case
of Moriwaki's construction of IndHilb-valued factorization homology
for conformally flat manifolds \cite{Mor26}.  Moriwaki proves that
for the Heisenberg algebra (and more generally for free-field
theories), the factorization homology on a Riemann surface $\Sigma_g$
takes values in the category $\mathrm{IndHilb}$ of ind-Hilbert
spaces, with the Bergman kernel providing the propagator and the
Fredholm determinant providing the partition function.

The present derivation recovers Moriwaki's result from the sewing
axioms and the Fock-space trace formula, without passing through
the full factorization homology machinery.  The advantage of the
sewing approach is its directness: the Fredholm determinant arises
from a one-line computation (the bosonic trace formula), while the
factorization homology approach requires the full construction of
the factorization algebra on $\Sigma_g$ and its global sections.
\end{remark}

\begin{remark}[Genus-$g$ generalization and the $\hat{A}$-genus]
% label removed: rem:genus-g-ahat
\index{A-hat genus!partition function}
For the Heisenberg algebra at genus $g \ge 2$, the partition function
generalizes to
\[
Z_g(\cH_k) = \det(1 - T_{\Omega_g})^{-k},
\]
where $T_{\Omega_g}$ is the restriction operator on the genus-$g$
surface with period matrix $\Omega_g$.  The Fredholm determinant
depends on $\Omega_g$ through a Siegel modular form.

The free energy $F_g = -\log Z_g = k\, \Tr\log(1 - T_{\Omega_g})$
has the asymptotic expansion
$F_g \sim k\, \sum_{n=1}^\infty \Tr(T_{\Omega_g}^n)/n$,
whose leading term is $k\, \Tr(T_{\Omega_g}) = k \cdot
(\text{dim }H^0(\Sigma_g, \Omega^1))/(24) = k\,g/24$
(by the Riemann--Roch trace formula).

At all genera, the generating function of modular characteristics
$\sum_{g \ge 1} F_g(\cH_k)\, x^{2g}$ is computed by
Volume~I, Theorem~D: it equals $k \cdot (\hat{A}(ix) - 1)$,
where $\hat{A}$ is the $\hat{A}$-genus (the multiplicative
sequence for the Todd class).  The Fredholm determinant formula
$Z_g = \det(1-T_q)^{-k}$ is the genus-by-genus evaluation of
this universal formula at each point of $\mathcal{M}_g$.
\end{remark}
